% concatenated from derived-transversality.tex
\documentclass[reqno]{amsart}

\usepackage[margin=1.5in]{geometry}
\usepackage{palatino}
\usepackage{mathpazo}

\usepackage{amsmath, amssymb, amsthm, mathrsfs}
\usepackage{tikz-cd}
\usepackage{hyperref}
\usepackage{enumerate}

\newtheorem{theorem}{Theorem}[section]
\newtheorem{lemma}[theorem]{Lemma}
\newtheorem{proposition}[theorem]{Proposition}
\newtheorem{observation}[theorem]{Observation}
\newtheorem{corollary}[theorem]{Corollary}
\newtheorem*{theorem*}{Theorem}

\theoremstyle{definition}
\newtheorem{definition}[theorem]{Definition}
\newtheorem{setup}[theorem]{Setup}
\newtheorem{example}[theorem]{Example}

\theoremstyle{remark}
\newtheorem{remark}[theorem]{Remark}

\newcommand{\LL}{\mathbb{L}}
\newcommand{\TT}{\mathbb{T}}
\newcommand{\Gm}{\mathbb{G}_m}
\newcommand{\cO}{\mathcal{O}}
\newcommand{\cL}{\mathcal{L}}
\newcommand{\cK}{\mathcal{K}}
\newcommand{\K}{\mathbb{K}}
\newcommand{\DR}{\mathbf{DR}}
\newcommand{\Spec}{\mathrm{Spec\,}}
\newcommand{\dArt}{\mathbf{dArt}}
\newcommand{\Acl}{\mathcal{A}^{2, \mathrm{cl}}}
\newcommand{\bfem}[1]{\textbf{\textit{#1}}}
\newcommand{\R}{{\mathbb R}}

\newcommand{\Arxiv}[1]{\href{http://arxiv.org/abs/#1}{#1}}

\begin{document}

\title[Transversality and Derived Legendrian Intersections]{Equivariant Quotients of derived symplectic spaces and \\ 
Legendrian Intersection Theorem}
\author{Efe \.{I}zbudak \and Kadri \.{I}lker Berktav}
\date{\today}

\begin{abstract}
The classical transversality lemma of contact geometry constructs a contact structure on a hypersurface transverse to a  Liouville vector field using point-set topology and local flows. This paper translates the classical transversality lemma into the context of derived algebraic geometry and provides the derived Legendrian intersection theorem, along with various applications to moduli theory. 
  
In brief,  we first prove that taking the quotient of a derived symplectic space descends the symplectic data to a contact structure, avoiding a transverse hypersurface, where  the fundamental vector field of a weight 1 $\Gm$-action, in the derived setting, replaces the classical Liouville vector field.  Secondly, the derived Legendrian intersection theorem is proven using base change, an $\infty$-categorical descent cube, and $\Gm$-equivariant lifts along the symplectification projection.
  
As applications of the main results, we first examine the derived geometry of the discriminant loci of 1-jet bundles and show that these loci carry a $(-1)$-shifted contact structure. In addition, we show that our results apply to certain moduli problems, including projective Higgs bundles, $\ell$-adic local systems, and Lie 2-groups, and we provide further examples of contact derived moduli stacks.
\end{abstract}

\maketitle
\tableofcontents

\section{Introduction and Summary}

A fundamental lemma in classical differential geometry \cite[Lemma 1.4.5]{Geiges} states that a Liouville vector field on a symplectic
manifold induces a contact structure on any transverse hypersurface.
This classical formulation  relies on smooth local flows, differential forms,
and point-set transversality \cite[Section
1.4]{Geiges}. These methods do not apply to moduli
problems in algebraic geometry, where spaces often have
singularities, nilpotence, and non-trivial automorphism groups. Similar drawbacks emerge for the intersection problems as well.
To address these issues,  nowadays folklore suggests, for instance, the use of shifted geometric (e.g., symplectic) structures, equivariant group actions, exact triangles, and derived cotangent complexes, or in short, the framework of Derived Algebraic Geometry.

It is well known that DAG essentially provides a new setting for dealing with non-generic situations in geometry, such as non-transversal intersections and “bad" quotients. It also offers generalized versions of certain familiar geometric structures with various outcomes. That approach is what we intend to adapt in this paper.
\vspace{0.05in}

\paragraph{\bfem{Conventions and notations}.} Throughout the paper, $ \mathbb{K} $ will be an algebraically closed field of characteristic zero. All cdgas will be graded in nonpositive degrees and  over $\mathbb{K}.$ We always consider $\K$-schemes/stacks, and we assume that all classical $ \K $-schemes are \textit{locally of finite type}, and that all derived $ \K $-schemes/stacks $ {X} $ are  \textit{locally finitely presented.}
\vspace{0.1in}

\paragraph{\bfem{Our results.}} We now summarize our main results, leaving full explanations and
definitions to the main text. Our first result shows that the classical transversality lemma of contact geometry can be reformulated in the derived context. More specifically, using the geometry of $\Gm$-torsors
and symplectification in the sense of \cite[Section 4.2]{Berktav1}, \cite[Theorem
4.1]{Berktav2}, we prove the following result.

\begin{theorem}[Derived Analogue of Classical Transversality]  \label{thm:A}
Given an $n$-shifted symplectic derived Artin stack $(\widetilde{X},
  \omega_{\widetilde{X}})$
  equipped with a $\Gm$-action of weight 1,
  the stack quotient $[\widetilde{X} / \Gm]$ inherits an
  $n$-shifted contact structure (cf. Theorem \ref{thm:derived transversality}).
\end{theorem}

This theorem, in fact, formalizes \textit{derived symplectification}. Taking the
quotient of a homogenous symplectic space descends the symplectic
data to a contact structure, avoiding
a transverse hypersurface. In the derived setting, the Liouville vector field corresponds to the
fundamental vector field of a weight 1 $\Gm$-action, and the
contact manifold corresponds to the stack quotient. Classical coordinate
arguments are replaced by stable $\infty$-categorical operations.
This relies on generalized quotient groupoids,
shifted symplectic forms \cite[Definition 1.18]{PTVV},
and exact triangles of perfect complexes.

The second main result of this paper concerns the \textit{Legendrian intersection problem}, which has no direct classical counterpart. In brief, the theory above extends to Legendrian morphisms and gives rise to an affirmative answer to the Legendrian intersection problem in the context of derived geometry. 
The central idea is that derived contact geometry links Legendrian structures to Lagrangian structures via $\Gm$-equivariant lifts along the symplectification projection.
Lifting intersection data to the symplectic level yields the following result.

\begin{theorem}[Legendrian Intersection Theorem] \label{thm:B}
Intersection of two Legendrians $f_1 \colon L_1 \to X$ and $f_2 \colon L_2 \to X$ in an $n$-shifted contact derived Artin stack admits 
 an $(n-1)$-shifted  contact structure (cf. Theorem \ref{thm:Leg intersection}).
\end{theorem}
It is important to note that this lifting property translates intersection problems in contact moduli to equivariant intersection problems in symplectic moduli. This establishes the desired shifted contact structure on the derived intersection of two Legendrian stacks.
As an interesting consequence of Theorem \ref{thm:B}, we then obtain:

\begin{corollary}\label{cor:B}
The \bfem{derived discriminant locus}
$ \Delta\mathrm{loc}(f):= L_F \times_{J^1(L)}^h L_0$ of a regular function $f$ on a smooth $\K$-scheme $L$ admits a $(-1)$-shifted contact structure, where  $F:p \mapsto (f(p),d_{dR}f_p)$ defines a Legendrian embedding $L_F$ of $L$ into the 1-jet space $(J^1(L), \xi_{jet})$ and $L_0$ denotes the Legendrian given by the zero-section morphism
    (cf. Example \ref{example: ddislocus}).
\end{corollary}

Last but not least, we also provide several applications of Theorem \ref{thm:A} and Theorem \ref{thm:B} to moduli theory and give examples of derived stacks with a shifted contact structure.

\begin{corollary} \label{cor:C}
    \begin{enumerate} \itemsep=5pt
        \item Let $X$ be a smooth proper variety of dimension $d$ over $\K$ and $G$ be a reductive algebraic group. The derived moduli stack $P\mathrm{Higgs}_G(X)$ of projective Higgs bundles admits a canonical $(1-d)$-shifted contact structure (cf. Corollary \ref{cor:projective_higgs}). 
        
        Restricting to a smooth proper curve $C$, the derived intersection of projectivized irreducible components of the global nilpotent cone within the $0$-shifted ambient contact stack $P\mathrm{Higgs}_G(C)$ yields a strictly derived moduli space equipped with a canonical $(-1)$-shifted contact structure (cf. Corollary \ref{cor:nilpotent_intersection}).
        \item Let $LocSys_{\ell,n}^{fr}(X)$ denote the \textit{derived moduli stack of framed geometric $\ell$-adic local systems on the base extension $X_{\bar{\mathbb{F}}_q}$}. Then the derived contact mapping torus $$\left[(LocSys_{\ell,n}^{fr}(X) \times \mathbb{A}^1[2-2d]) / \mathbb{Z}\right]$$ inherits a $(2-2d)$-shifted contact structure (cf. Corollary \ref{cor: l-adic local sys}).        \item The derived stack quotient $\mathbb{P}(T^*[2]BG) := [T^*[2]BG^\circ / \Gm]$   carries a canonical 2-shifted contact structure (cf. Corollary \ref{cor: Lie 2-group}).
        \item For a closed subgroup $H \subset G$, denote by $\mathcal{N}^*_{BH/BG}[2]$ the 2-shifted conormal stack. Then  the morphism $\mathbb{P}(\mathcal{N}^*_{BH/BG}[2]) \to \mathbb{P}(T^*[2]BG)$ is Legendrian. Furthermore, the derived fiber product of two such Legendrians associated with subgroups $H_1$ and $H_2$ carries a canonical 1-shifted contact structure (cf. Corollary \ref{cor: intersection of prequan}).
    \end{enumerate}
\end{corollary}

\vspace{0.1in}
\paragraph{\bfem{Organization of the paper.}}  Sections \ref{sec:stable_cats} through
\ref{sec:structures} provide background material on stable
$\infty$-categories, derived commutative algebra, cotangent
complexes, and shifted symplectic and contact structures. Section
\ref{sec:classical} overviews the classical setup of transversality
in contact geometry. Section
\ref{sec:setup} formulates the derived setup and the first main theorem of this paper (cf. Theorem \ref{thm:derived transversality}). In
Section \ref{sec:proof}, we give the proof of Theorem \ref{thm:A} and
analyze the failure of classical transversality.  Section
\ref{sec:legendrian} establishes  Theorem \ref{thm:B} and Corollary \ref{cor:B}. Finally, in Section \ref{sec:moduli}, we elaborate on the content of Corollary \ref{cor:C} and give the proof of each statement.

\section{Recollection}

\subsection{Elements of Stable \texorpdfstring{$\infty$-}-Category Theory} \label{sec:stable_cats}

In derived geometry, sets and sheaves are
replaced by stable $\infty$-categories. Let $X$ be a topological space or a scheme.

\begin{definition}[Stable $\infty$-Categories]
  Following Lurie \cite[Section 1.1.1]{LurieHA}, a \bfem{stable
  $\infty$-category} is an
  $\infty$-category $\mathcal{C}$ possessing a zero object, as well
  as all finite limits and colimits, such that a
  commutative square in $\mathcal{C}$ is a pullback if and only if it
  is a pushout.
  
  Here, the suspension functor $\Sigma: \mathcal{C}
  \xrightarrow{\sim} \mathcal{C}$
  is an equivalence, allowing every morphism to be
  completed to a fiber sequence.
\end{definition}

\begin{definition}[The Derived $\infty$-Category]
  For a derived stack $X$, the $\infty$-category of quasi-coherent
  sheaves is denoted $L_{qcoh}(X)$ \cite[Section 1.1]{PTVV}. Its
  homotopy category is the quasi-coherent derived category
  $D_{qcoh}(X)$. $L_{qcoh}(X)$ is a stable, symmetric monoidal
  $\infty$-category possessing internal mapping objects
  $\mathbb{R}\underline{\mathcal{H}om}(E, F)$.
\end{definition}

This localization
resolves the classical failure of flat base change, permitting the
global computation of
derived functors.

\begin{definition}[Internal Mapping Objects and Mapping Spaces]
  In a stable, symmetric monoidal $\infty$-category $\mathcal{C}$
  equipped with a closed structure, the internal mapping object
  $\mathbb{R}\underline{\mathcal{H}om}(E, F)$ resides within
  $\mathcal{C}$. The \bfem{mapping space}
  $\mathrm{Map}_{\mathcal{C}}(E, F)$ is a Kan complex in the
  $\infty$-category of spaces $\mathcal{S}$. By the stable
  $\infty$-categorical Yoneda lemma, morphisms in $\mathcal{C}$ are
  identified with global sections of the internal mapping object.
\end{definition}

\begin{definition}[Homotopy Fibers and Fiber Sequences]
  Let $f: E \to F$ be a morphism in the stable $\infty$-category $L_{qcoh}(X)$.
  The \bfem{homotopy fiber} $\mathrm{fib}(f)$ is defined
  as the limit
  (homotopy pullback) of the diagram $E \xrightarrow{f} F \leftarrow
  0$. 
  
  This construction
  gives a \bfem{fiber sequence} $\mathrm{fib}(f) \to E
  \xrightarrow{f} F$, which
  serves as the derived replacement for a short exact sequence \cite[Section 1.1.2]{LurieHA}.
\end{definition}

In a stable $\infty$-category, fiber sequences and cofiber sequences
coincide up to a shift. This endows the category with a
triangulated structure, allowing morphisms to be completed to
an exact sequence. The homotopy cofiber provides
a method for computing homological obstructions.

\begin{remark}
  Fiber sequences replace short exact sequences in derived geometry.
  Classically, the kernel or cokernel of a map of vector bundles on a
  singular space might fail to be a vector bundle. The homotopy cofiber
  replaces the cokernel, and the homotopy fiber
  replaces the kernel. Unlike kernels and cokernels, homotopy
  fibers and cofibers preserve
  dualizability.
\end{remark}

\begin{lemma}[Pasting Law for Pullbacks]
  \label{lem:pullback_pasting}
  Let $\mathcal{C}$ be an $\infty$-category and suppose we are given a diagram $\Delta^1 \times \Delta^2 \to \mathcal{C}$ depicted as
  \begin{align*}
    \begin{tikzcd}[ampersand replacement=\&]
      A \arrow[r] \arrow[d] \& B \arrow[r] \arrow[d] \& C \arrow[d] \\
      X \arrow[r] \& Y \arrow[r] \& Z
    \end{tikzcd}
  \end{align*}
  If the right square is a pullback in $\mathcal{C}$, then the left square is a pullback if and only if the outer rectangle is a pullback.
\end{lemma}
\begin{proof}
  This follows by applying the pasting law for pushouts \cite[Lemma 4.4.2.1]{LurieHTT} to the opposite $\infty$-category $\mathcal{C}^{op}$. A pullback square in $\mathcal{C}$ is a pushout square in $\mathcal{C}^{op}$. The diagram $\Delta^1 \times \Delta^2 \to \mathcal{C}$ corresponds to a diagram $(\Delta^1 \times \Delta^2)^{op} \simeq \Delta^1 \times \Delta^2 \to \mathcal{C}^{op}$. Applying \cite[Lemma 4.4.2.1]{LurieHTT} in $\mathcal{C}^{op}$ finishes the proof.
\end{proof}

\begin{lemma}[Exact Triangles of Fibers and the Pasting Law]
  \label{lem:pasting_law}
  Let $\mathcal{C}$ be a stable $\infty$-category. Given composable
  morphisms $X \xrightarrow{f} Y \xrightarrow{g} Z$ in $\mathcal{C}$,
  the homotopy fibers assemble into a fiber sequence
  \begin{align*}
    \mathrm{fib}(f) \to \mathrm{fib}(g \circ f) \to \mathrm{fib}(g).
  \end{align*}
\end{lemma}
\begin{proof}
  Since $\mathcal{C}$ is a stable $\infty$-category, it admits all finite limits \cite[Proposition 1.1.3.4]{LurieHA}. 
  
  \noindent The composition $g \circ f$ determines a  morphism of cospans from $(X \xrightarrow{g \circ f} Z \leftarrow 0)$ to $$(Y \xrightarrow{g} Z \leftarrow 0).$$ 
  
  The functoriality of limits extends this to a diagram $\Delta^1 \times \Delta^2 \to \mathcal{C}$ containing two adjacent squares
  \begin{align*}
    \begin{tikzcd}[ampersand replacement=\&]
      \mathrm{fib}(g \circ f) \arrow[r] \arrow[d] \& \mathrm{fib}(g) \arrow[r] \arrow[d] \& 0 \arrow[d] \\
      X \arrow[r, "f"] \& Y \arrow[r, "g"] \& Z
    \end{tikzcd}
  \end{align*}
  The right square is a pullback by the definition of the homotopy fiber $\mathrm{fib}(g)$ \cite[Definition 1.1.1.6]{LurieHA}.
  The outer rectangle is a pullback by the definition of the homotopy fiber $\mathrm{fib}(g \circ f)$.
  
  Lemma \ref{lem:pullback_pasting} implies the left square
  \begin{align*}
    \begin{tikzcd}[ampersand replacement=\&]
      \mathrm{fib}(g \circ f) \arrow[r] \arrow[d] \& \mathrm{fib}(g) \arrow[d] \\
      X \arrow[r, "f"] \& Y
    \end{tikzcd}
  \end{align*}
  is a pullback square.
  We then extend the diagram upwards by defining $F$ as the homotopy fiber of the morphism $\mathrm{fib}(g \circ f) \to \mathrm{fib}(g)$, obtaining the diagram
  \begin{align*}
    \begin{tikzcd}[ampersand replacement=\&]
      F \arrow[r] \arrow[d] \& 0 \arrow[d] \\
      \mathrm{fib}(g \circ f) \arrow[r] \arrow[d] \& \mathrm{fib}(g) \arrow[d] \\
      X \arrow[r, "f"] \& Y
    \end{tikzcd}
  \end{align*}
  The top square is a pullback by the definition of $F$. The middle square is a pullback as demonstrated above. 
  
  Applying Lemma \ref{lem:pullback_pasting} again to this vertical arrangement, the outer rectangle
  \begin{align*}
    \begin{tikzcd}[ampersand replacement=\&]
      F \arrow[r] \arrow[d] \& 0 \arrow[d] \\
      X \arrow[r, "f"] \& Y
    \end{tikzcd}
  \end{align*}
  is a pullback square. By the definition of the homotopy fiber, we obtain an equivalence $F \simeq \mathrm{fib}(f)$.

  The definition of $F$ as the fiber of $\mathrm{fib}(g \circ f) \to \mathrm{fib}(g)$ provides the fiber sequence
  \begin{align*}
    F \to \mathrm{fib}(g \circ f) \to \mathrm{fib}(g).
  \end{align*}
  Since $F \simeq \mathrm{fib}(f)$, we obtain the fiber sequence
  \begin{align*}
    \mathrm{fib}(f) \to \mathrm{fib}(g \circ f) \to \mathrm{fib}(g).
  \end{align*}
\end{proof}

\begin{definition}[Perfect Complexes]
  A \bfem{perfect complex} on a derived stack $X$ is
  defined as a \textbf{dualizable object} in the symmetric monoidal
  $\infty$-category $L_{qcoh}(X)$ \cite[Section 1.1]{PTVV}. This
  means there exists an object $E^\vee :=
  \mathbb{R}\underline{\mathcal{H}om}(E, \mathcal{O}_X)$ such that
  the biduality map $E \to (E^\vee)^\vee$ is an equivalence.
\end{definition}

Perfect complexes represent the dualizable objects in the derived
category of quasi-coherent sheaves. They provide the geometric
finiteness condition required in shifted symplectic geometry
\cite[Section 1.1]{Calaque}. Without this
condition, the biduality map fails to be an equivalence, precluding the
definition of a non-degenerate symplectic form.

\subsection{Derived Commutative Algebras and Intersections} \label{sec:derived_comm}

\begin{definition}[Derived Tensor Product]
  Let $M, N \in L_{qcoh}(X)$. The \bfem{derived tensor product} $M
  \otimes^{\mathbb{L}}_{\cO_X} N$
  is the canonical symmetric monoidal structure on the stable
  $\infty$-category $L_{qcoh}(X)$ \cite[Section 4.4.2]{LurieHA}.
  Its cohomology group in degree $-i$ computes the higher derived functor
  $\mathrm{Tor}_i^{\cO_X}(M, N)$.
\end{definition}

The derived tensor product is the algebraic
operation for derived intersection theory. Computing intersections
via derived tensor products preserves geometric intersection
multiplicities \cite[Section 1.1]{PTVV}. This
extends the intersection theory of Serre by
embedding intersection multiplicities into the structure sheaf.

\begin{definition}[Commutative Differential Graded Algebra]
  A \bfem{commutative differential graded algebra} (cdga) over $\K$
  is a graded $\K$-algebra $A = \bigoplus_{i \le 0} A^i$ concentrated
  in non-positive cohomological degrees, equipped with a differential
  $d: A^i \to
  A^{i+1}$ satisfying $d^2 = 0$ and the graded Leibniz rule.
  The $\infty$-category of such algebras is denoted by
  $\mathrm{cdga}_\K^{\leq 0}$.
\end{definition}
These algebras generalize commutative rings by
encoding relations and syzygies into the differential
structure. The study of affine derived schemes is formally dual to
the study of these commutative differential graded algebras,
providing an algebraic framework for derived geometry.

\begin{definition}[Homotopy Pullback]
  Let $X \to Z$ and $Y \to Z$ be morphisms of derived spaces. The
  \bfem{homotopy pullback} or derived intersection $X
  \times_Z^{\mathbb{R}} Y$ is the limit completing the
  pullback diagram in the $\infty$-category of derived spaces. 
  
  Affinely, if
  $X = \Spec(A)$, $Y = \Spec(B)$, and $Z = \Spec(C)$, the derived
  intersection corresponds to the pushout $A
  \otimes_C^{\mathbb{L}} B$ computed within
  $\mathrm{cdga}_\K^{\leq 0}$
  \cite[Section 2.2.2]{HAG-II}.
\end{definition}

The universal property of the homotopy pullback ensures fiber
products of derived stacks exist. This framework resolves classical
pathologies where fiber
products fail to be smooth or exhibit unexpected dimensions. The
resulting space is a derived scheme
capturing the intersection geometry.

\begin{remark}
  Let $X$ and $Y$ be closed subschemes of a scheme $Z$. The classical intersection $X \times_Z Y$ has the structure sheaf $\mathcal{O}_X \otimes_{\mathcal{O}_Z} \mathcal{O}_Y$. If $X$ and $Y$ do not intersect transversely, the tensor product $\mathcal{O}_X \otimes_{\mathcal{O}_Z} \mathcal{O}_Y$ discards homological data. The resulting scheme lacks the necessary structure to determine the intersection multiplicity.
\end{remark}

\begin{observation}
  \label{prop:cdga_necessity}
  Let $\mathbb{K}$ be a commutative ring. The derived fiber product of the origin $\mathrm{Spec}(\mathbb{K})$ with itself over the affine line $\mathbb{A}^1_{\mathbb{K}}$ is represented by the commutative differential graded algebra $\mathbb{K}[\epsilon]/(\epsilon^2)$ with $\epsilon$ in degree $-1$. The homology of this algebra recovers the intersection multiplicity.
\end{observation}
\begin{proof}
  The classical fiber product $\mathrm{Spec}(\mathbb{K}) \times_{\mathrm{Spec}(\mathbb{K}[x])} \mathrm{Spec}(\mathbb{K})$ is computed by the pushout in the category of commutative rings
  \begin{align*}
    \mathbb{K} \otimes_{\mathbb{K}[x]} \mathbb{K} \cong \mathbb{K}.
  \end{align*}
  This pushout evaluates to $\mathbb{K}$ and discards the non-transverse intersection multiplicity.

  The derived intersection is governed by the derived tensor product $\mathbb{K} \otimes_{\mathbb{K}[x]}^{\mathbb{L}} \mathbb{K}$ in the $\infty$-category of commutative differential graded algebras. We compute this derived tensor product by replacing $\mathbb{K}$ with the Koszul resolution $K(\mathbb{K}[x], x)$ over $\mathbb{K}[x]$
  \begin{align*}
    0 \to \mathbb{K}[x] \xrightarrow{x} \mathbb{K}[x] \to 0
  \end{align*}
  where the terms are concentrated in cohomological degrees $-1$ and $0$. Applying the functor $- \otimes_{\mathbb{K}[x]} \mathbb{K}$ to the Koszul resolution we obtain the chain complex
  \begin{align*}
    0 \to \mathbb{K} \xrightarrow{0} \mathbb{K} \to 0.
  \end{align*}
  The homology of this complex computes the higher derived functors
  \begin{align*}
    \mathrm{Tor}_0^{\mathbb{K}[x]}(\mathbb{K}, \mathbb{K}) &\cong \mathbb{K} \\
    \mathrm{Tor}_1^{\mathbb{K}[x]}(\mathbb{K}, \mathbb{K}) &\cong \mathbb{K}.
  \end{align*}
  The derived intersection possesses a structure sheaf represented by the commutative differential graded algebra $\mathbb{K}[\epsilon]/(\epsilon^2)$ with generator $\epsilon$ in degree $-1$ and differential $d(\epsilon) = 0$. The $\infty$-categorical theory preserves the homological data of the non-transverse intersection \cite[Section 2]{ToenDAG}.
\end{proof}

\subsection{From Classical Differentials to Derived Stacks} \label{sec:diffs}

\begin{definition}[K\"{a}hler Differentials]
  Let $A$ be a commutative $R$-algebra. The module of
  \bfem{K\"{a}hler differentials} $\Omega^1_{A/R}$ is the
  $A$-module generated by symbols $da$ subject to the
  Leibniz rule. It is not homotopy invariant for singular spaces,
  making it inadequate for
  derived moduli theory \cite[Section 1.3.1]{CPTVV}.
\end{definition}

The module of K\"{a}hler differentials is locally free for
smooth schemes, but behaves poorly for singularities.
Its lack of exactness requires a derived replacement for moduli
problems \cite[Section 1.1]{Calaque}. The
failure of the cotangent functor to be left exact means the
classical sheaf cannot distinguish between singular
structures with identical tangent spaces.

In what follows, we present the underlying (higher) spaces of interest in this paper and outline their essential properties, along with several useful remarks.

\begin{definition}[Affine Derived Schemes]
  The $\infty$-category of \bfem{affine derived schemes} over $\K$,
  denoted $\mathbf{dAff}_\K$, is defined as the opposite
  $\infty$-category $(\mathrm{cdga}_\K^{\leq 0})^{\mathrm{op}}$. For
  any $A \in \mathrm{cdga}_\K^{\leq 0}$, the corresponding affine
  derived scheme is denoted $\Spec(A)$ \cite[Definition 2.2.2.4]{HAG-II}.
\end{definition}

\begin{definition}[Derived Stacks]
  A \bfem{derived stack} over $\K$ is a functor $X \colon
  \mathrm{cdga}_\K^{\leq 0} \to \mathcal{S}$ taking values in the
  $\infty$-category of spaces (Kan complexes), such that $X$
  satisfies exact hyperdescent with respect to the \'{e}tale
  topology. The $\infty$-category of derived stacks, denoted
  $\mathbf{dSt}_\K$, is constructed as the left Bousfield
  localization of the $\infty$-category of presheaves
  $\mathrm{Fun}(\mathrm{cdga}_\K^{\leq 0}, \mathcal{S})$ along
  \'{e}tale hypercovers \cite[Definition 2.2.2.14]{HAG-II}.
\end{definition}
In other words, a derived stack $ {X} $ is  a  functor  $$ {X}:  \mathrm{cdga}_{\K}^{\leq 0} \rightarrow  \mathcal{S}, \ \ A\mapsto X(A) \simeq \mathrm{Map}_{\mathbf{dSt}_\K}(\Spec(A), {X}),$$ satisfying a descent condition.
For more details, we refer to \cite{HAG-II}. 
It should then be noted that any affine derived $\mathbb{K}$-scheme $X$ can be corepresented by a cdga $A$, where $\Spec(A)$ is given as the functor \[ \Spec (A): B\in \mathrm{cdga}_\K^{\leq 0}  \longmapsto Hom_{\mathrm{cdga}_{\mathbb{K}}^{\leq 0}} (A,B) \in \mathcal{S}.\]

The concept of a derived stack may sometimes be too wild for some purposes because a typical $A$-point $p\in X(A)$ is represented by a general morphism $p: \Spec A \to X$ in $\mathbf{dSt}_\K$. When we ask for additional conditions on those morphisms, such as being a Zariski open embedding or being smooth (of some relative dimension), the resulting space may admit a tractable description in terms of affine covers, leading to the following special kind of derived stacks.

\begin{definition}[Derived Schemes]
  \label{def:derived_schemes}
  An object $X$ in the $\infty$-category of derived stacks $\mathbf{dSt}_\K$ is a \bfem{derived scheme} over $\K$ if it can be covered by Zariski open affine derived schemes $Y \subset X$. The $\infty$-category of derived schemes, denoted $\mathbf{dSch}_\K$, is the full subcategory of $\mathbf{dSt}_\K$ spanned by the derived schemes \cite[Definition 2.2]{Berktav2}.
\end{definition}

\begin{definition}[Smooth Morphisms of Derived Stacks]
  A morphism of derived stacks $f: X \to Y$ is \bfem{smooth} if it
  is locally of finite presentation and its relative cotangent
  complex $\LL_{X/Y}$ is a perfect complex of Tor-amplitude $[0, 0]$,
  meaning it acts homotopically as a locally free sheaf
  \cite[Section 2.2.5]{HAG-II}.
\end{definition}

\begin{definition}[Derived Artin Stacks]
  A \bfem{derived Artin stack} is a derived stack
  which is $m$-geometric for some integer $m$, for the \'{e}tale
  topology, and for the class of smooth maps \cite[Definition
  1.3.1]{PTVV}. 
\end{definition}
As mentioned before, all derived $ \K $-schemes/stacks $ {X} $ are assumed to be  \textit{locally finitely presented.} Let us also review key properties of cotangent and tangent complexes associated to these higher spaces.
\begin{definition}[Trivial Square-Zero Extensions and Derived Derivations]
  Let $X$ be a derived stack and $M \in L_{qcoh}(X)$. The
  \bfem{trivial square-zero extension} of $X$ by $M$, denoted
  $X[M]$, is defined as the relative derived spectrum
  $\mathbb{R}\mathbf{Spec}_X(\cO_X \oplus M)$, where $M$ acts as a
  square-zero ideal. The space of \bfem{derived derivations}
  $\mathrm{Der}(X, M)$ is the mapping space of sections
  $\mathrm{Map}_{\dArt_{\K // X}}(X, X[M])$ \cite[Section 1.3.1]{CPTVV}.
\end{definition}

\begin{definition}[Cotangent and Tangent Complexes]
  For a derived Artin stack $X$, assumed to be locally of finite
  presentation over $\K$, the \bfem{cotangent complex}
  $\LL_{X/\K} \in L_{qcoh}(X)$ is the complex defined by
  corepresenting derived derivations. It satisfies the adjunction
  $\mathrm{Map}_{L_{qcoh}(X)}(\LL_{X/\K}, M) \simeq
  \mathrm{Map}_{\dArt_{\K // X}}(X, X[M])$ \cite[Section
  1.3.1]{CPTVV}. Since $X$ is locally of finite presentation,
  $\LL_{X/\K}$ is a perfect complex and thus dualizable. Its dual is the
  \bfem{tangent complex} $\TT_{X/\K} \simeq
  \mathbb{R}\underline{\mathcal{H}om}(\LL_{X/\K}, \cO_X)$, which
  classifies derived vector fields.
\end{definition}

The cotangent complex provides a derived replacement for
K\"{a}hler differentials. By extending classical differentials
into a perfect complex, it captures the obstruction theory of
the space \cite[Section 1.1]{Calaque}. The higher cohomology groups
contain the obstruction classes for deformations, while
the degree zero part recovers the classical differentials.

\begin{observation}
  \label{prop:cotangent_necessity}
  Let $A = R/I$ be a quotient of a commutative ring $R$. The module of K\"{a}hler differentials $\Omega^1_A$ gives a right-exact conormal sequence. The derived cotangent complex $\mathbb{L}_A$ lifts this to an exact triangle in the derived category of $A$-modules.
\end{observation}
\begin{proof}
  The classical module of K\"{a}hler differentials fits into the conormal sequence
  \begin{align*}
    I/I^2 \to \Omega^1_R \otimes_R A \to \Omega^1_A \to 0.
  \end{align*}
  This sequence is right-exact. The left morphism maps a class $\overline{f} \in I/I^2$ to $df \otimes 1$. If $A$ is singular, this morphism is generally not injective. The kernel corresponds to the relations among the generators of $I$.

  The derived cotangent complex $\mathbb{L}_A$ extends this sequence. It is constructed by taking a simplicial resolution of $A$ by free commutative rings over $R$ and applying the K\"{a}hler differentials functor degree-wise. The ring homomorphism $R \to A$ induces the transitivity sequence, which is an exact triangle in the derived category of $A$-modules
  \begin{align*}
    \mathbb{L}_R \otimes_R^{\mathbb{L}} A \to \mathbb{L}_A \to \mathbb{L}_{A/R}.
  \end{align*}
  Assume $A$ is a global complete intersection. The relative cotangent complex $\mathbb{L}_{A/R}$ is quasi-isomorphic to the shifted conormal module $(I/I^2)[1]$ \cite[Section 1.1]{Calaque}. We apply the shift functor $[-1]$ to rotate the exact triangle and obtain
  \begin{align*}
    \mathbb{L}_{A/R}[-1] \to \mathbb{L}_R \otimes_R^{\mathbb{L}} A \to \mathbb{L}_A.
  \end{align*}
  Substituting the quasi-isomorphism $\mathbb{L}_{A/R}[-1] \simeq I/I^2$ results in the exact triangle
  \begin{align*}
    I/I^2 \to \mathbb{L}_R \otimes_R^{\mathbb{L}} A \to \mathbb{L}_A.
  \end{align*}
  Applying the homology functor to this triangle produces a long exact sequence. The zeroth homology modules are $H_0(\mathbb{L}_R \otimes_R^{\mathbb{L}} A) \cong \Omega^1_R \otimes_R A$ and $H_0(\mathbb{L}_A) \cong \Omega^1_A$. The long exact sequence terminates as
  \begin{align*}
    H_1(\mathbb{L}_A) \to I/I^2 \to \Omega^1_R \otimes_R A \to \Omega^1_A \to 0.
  \end{align*}
  This sequence recovers the classical conormal sequence and extends it to the left. The module $H_1(\mathbb{L}_A)$ measures the failure of injectivity. 
  
\end{proof}
We now give several definitions, which will be central for the rest of the paper.
\begin{definition}
  \label{def:derived_group_scheme}
  Let $S$ be a base derived scheme. A \bfem{derived group scheme}, or simply a \bfem{group scheme} over $S$ is a functor $G \colon \mathrm{cdga}_{\mathbb{K}}^{\leq 0} \to \mathcal{S}$ that factors through $\mathrm{Grp}(\mathcal{S})$, which denotes the $\infty$-category of group objects in $\mathcal{S}$, such that the underlying functor to $\mathcal{S}$ is representable by a derived scheme $X$ over $S$. 
  
  If $X$ is affine, then $G$ is a \bfem{derived affine group scheme}. If the structure map $X\to S$ is smooth then $G$ is a \bfem{smooth derived group scheme}. \cite[adapted from App. A.2]{Iwanari}.
\end{definition}

We will usually confuse the functor $G$ with the derived scheme $X$ in the above definition of a derived group scheme since by the derived Yoneda lemma the $\infty$-category of derived schemes embeds fully faithfully into the $\infty$-category of presheaves of spaces. 

\begin{definition}
  \label{def:multiplicative_group}
  Let $S = \Spec(\K)$. The \bfem{multiplicative group scheme} $\mathbb{G}_m$ over $S$ is the smooth derived affine group scheme represented by the commutative differential graded algebra $\K[t, t^{-1}]$ concentrated in degree zero. 
  
  The group scheme structure is determined functorially. For any $A \in \mathrm{cdga}_{\K}^{\leq 0}$, the space of points $\mathbb{G}_m(A)$ is the group $H^0(A)^\times$ of invertible elements in $H^0(A)$ \cite[adapted from Section 1.2]{Waterhouse}.
\end{definition}

\begin{definition}
  \label{def:quotient_groupoid}
  Let $G$ be a smooth group scheme over a base derived scheme $S$. Assume that $G$ acts on a derived scheme $X$ over $S$. The \bfem{quotient groupoid} $B_\bullet(X, G)$ is the simplicial object with degree $n$ terms $X \times_S G^n$. 
  
  The face maps $d_i \colon X \times_S G^n \to X \times_S G^{n-1}$ and degeneracy maps $s_i \colon X \times_S G^n \to X \times_S G^{n+1}$ are defined on points $(x, g_1, \dots, g_n)$ by
  \begin{align*}
    d_0(x, g_1, \dots, g_n) &= (g_1 \cdot x, g_2, \dots, g_n) \\
    d_i(x, g_1, \dots, g_n) &= (x, g_1, \dots, g_i g_{i+1}, \dots, g_n) \quad 0 < i < n \\
    d_n(x, g_1, \dots, g_n) &= (x, g_1, \dots, g_{n-1}) \\
    s_i(x, g_1, \dots, g_n) &= (x, g_1, \dots, g_i, e, g_{i+1}, \dots, g_n) \quad 0 \leq i \leq n
  \end{align*}
  where $g_1 \cdot x$ denotes the action of $G$ on $X$ and $e$ is the identity of $G$. The object $B_\bullet(X, G)$ is a smooth groupoid of derived schemes over $S$.
\end{definition}

\begin{definition}
  \label{def:derived_quotient_stack}\cite[Section 3.3]{ToenDAG}
  Let $G$ be a smooth group scheme over a base derived scheme $S$ acting on a derived scheme $X$ over $S$. The \bfem{quotient stack} $[X/G]$ is the derived Artin stack $|B_\bullet(X, G)|$,
   defined as the colimit of the simplicial object $B_\bullet(X, G)$:
    \begin{equation*}
        \big[{X}/G\big]= \mathrm{colim} \bigg(X\mathrel{\substack{\textstyle\leftarrow\\[-0.1ex]
                \textstyle\leftarrow \\[-0.1ex]}} X \times_S G
        \mathrel{\substack{\textstyle\leftarrow\\[-0.1ex]
                \textstyle\leftarrow \\[-0.1ex]
                \textstyle\leftarrow\\[-0.3ex]}} X \times_S G^2
        \mathrel{\substack{\textstyle\leftarrow\\[-0.1ex]
                \textstyle\leftarrow \\[-0.1ex]
                \textstyle\leftarrow \\[-0.1ex]
                \textstyle\leftarrow \\[-0.3ex]}} 
        \cdot\cdot\cdot \bigg),
    \end{equation*}  where the maps are given by the action and projection.
\end{definition} 
\begin{definition}
  \label{def:classifying_stack}
  Let $G$ be a smooth group scheme over a base derived scheme $S$. The \bfem{classifying stack} $BG$ is the quotient stack $[S/G]$ defined by the trivial action of $G$ on $S$, \cite[Section 3.3]{ToenDAG}.
\end{definition}

The stack quotient integrates with the theory of perfect complexes, replacing the classical requirement of global slices.

\begin{observation}
  \label{prop:stack_necessity}
  Let $\mathbb{K}$ be a field and let $S = \mathrm{Spec}(\mathbb{K})$. Let $\mathbb{G}_m = \mathrm{Spec}(\mathbb{K}[t, t^{-1}])$ be the multiplicative group scheme over $S$. The orbit space $\mathbb{A}^1_S / \mathbb{G}_m$ in the category of ringed spaces is non-Hausdorff. The quotient stack $[\mathbb{A}^1_S / \mathbb{G}_m]$ restricts to the classifying stack $B\mathbb{G}_m$ at the origin and to the scheme $S$ on the open orbit.
\end{observation}
\begin{proof}
  Let the smooth group scheme $\mathbb{G}_m$ over $S$ act on the affine line $\mathbb{A}^1_S$ by multiplication. The action admits two orbits. The first orbit is the origin, which we identify with $S$. The second orbit is the open complement $U = \mathbb{A}^1_S \setminus S$. The $\mathbb{G}_m$-invariant open sets in $\mathbb{A}^1_S$ are the empty set, $U$, and $\mathbb{A}^1_S$. These invariant sets determine the quotient topology on the orbit space. The only open set containing the origin is the entire space $\mathbb{A}^1_S$, which necessarily contains the open orbit $U$. The origin and the open orbit cannot be separated by disjoint open sets. The orbit space is then non-Hausdorff.

  The quotient stack $[\mathbb{A}^1_S / \mathbb{G}_m]$ is the derived Artin stack $|B_\bullet(\mathbb{A}^1_S, \mathbb{G}_m)|$. We restrict this quotient stack to the invariant subschemes.

  Over the open orbit $U$, the action is simply transitive. The restricted quotient groupoid $B_\bullet(U, \mathbb{G}_m)$ is an equivalence relation. The quotient stack $|B_\bullet(U, \mathbb{G}_m)|$ is the scheme-theoretic quotient $S$.

  Over the origin, the inclusion $S \hookrightarrow \mathbb{A}^1_S$ is $\mathbb{G}_m$-equivariant with respect to the trivial action on $S$. The restriction of the quotient groupoid $B_\bullet(\mathbb{A}^1_S, \mathbb{G}_m)$ to the origin yields the smooth groupoid of derived schemes $B_\bullet(S, \mathbb{G}_m)$. The quotient stack $|B_\bullet(S, \mathbb{G}_m)|$ is the classifying stack $B\mathbb{G}_m$.   
  
\end{proof}
\vspace{1in}

\begin{definition}
  \label{def:equivariant_sheaves}
  Let $G$ be a smooth group scheme over a base derived scheme $S$ acting on a derived scheme $X$ over $S$. The stable $\infty$-category of \bfem{$G$-equivariant quasi-coherent sheaves} on $X$, denoted $L_{qcoh}^G(X)$, is defined as the limit of the cosimplicial $\infty$-category $L_{qcoh}(B_\bullet(X, G))$. This cosimplicial diagram is obtained by applying the contravariant functor $L_{qcoh}(-)$ to the quotient groupoid $B_\bullet(X, G)$.
\end{definition}
That is, $L_{qcoh}^G(X)$ is defined as the limit of the cosimplicial object $L_{qcoh}(B_\bullet(X, G))$
\begin{equation*}
    L_{qcoh}^G(X) = \lim \bigg(L_{qcoh}(X)\mathrel{\substack{\textstyle\rightarrow\\[-0.1ex]
            \textstyle\rightarrow \\[-0.1ex]}} L_{qcoh}(X \times_S G)
    \mathrel{\substack{\textstyle\rightarrow\\[-0.1ex]
            \textstyle\rightarrow \\[-0.1ex]
            \textstyle\rightarrow\\[-0.3ex]}} L_{qcoh}(X \times_S G^2)
    \mathrel{\substack{\textstyle\rightarrow\\[-0.1ex]
            \textstyle\rightarrow \\[-0.1ex]
            \textstyle\rightarrow \\[-0.1ex]
            \textstyle\rightarrow \\[-0.3ex]}} 
    \cdot\cdot\cdot \bigg)
\end{equation*}
where the maps are the pullbacks induced by the action and projection morphisms.

\begin{definition}
  \label{def:principal_torsor}
  Let $G$ be a derived group scheme over a base derived scheme $S$. A morphism of derived stacks $p \colon \widetilde{X} \to X$ equipped with a $G$-action on $\widetilde{X}$ over $X$ is a \bfem{principal $G$-torsor} if the natural morphism $[\widetilde{X}/G] \to X$ is an isomorphism in the homotopy category of derived stacks and the square
  \begin{equation*}
    \begin{tikzcd}
      \widetilde{X} \arrow[r] \arrow[d, "p"'] & S \arrow[d] \\
      X \arrow[r] & BG
    \end{tikzcd}
  \end{equation*}
  is homotopy cartesian, where $S \to BG$ is the canonical morphism from the base scheme to the classifying stack \cite[Definition 1.3.5.2]{TV2}.
\end{definition}

Note that the homotopy cartesian square condition is equivalent to the requirement that the shear map $G \times_S \widetilde{X} \to \widetilde{X} \times_X \widetilde{X}$ defined by $(g, x) \mapsto (gx, x)$ is an equivalence.

When $X\in \mathbf{dAff}_\K$,  say $X=\Spec A$ with $A$ a cdga,  a \bfem{$G$-torsor over $A$} is defined to be a $G$-module $F \in G\text{-Mod}$, together with a fibration of $G$-modules
\(
F \longrightarrow \Spec A,
\)
such that there exists an étale covering $A \longrightarrow B$ with the property that the object
\[
F \times_{\Spec A} \Spec B \longrightarrow \Spec B
\]
is equivalent over $\Spec B$ to
\(
G \times \Spec B \longrightarrow \Spec B
\) (where $G$ acts on itself by left translations). For details, we refer to \cite[Section 3.4]{TV2}.

\subsection{Derived Symplectic, Lagrangian, and Contact Structures}
\label{sec:structures}

\begin{definition}[Derived de Rham Complex and Closed Forms]
  For $A \in \mathrm{cdga}_\K^{\leq 0}$, the \bfem{derived de Rham
  algebra} $\DR(A/\K)$ is the graded mixed complex evaluated as
  $\mathrm{Sym}_A(\LL_{A/\K}[1])$ incorporating the standard Hodge
  filtration \cite[Section 1.4.1]{CPTVV}.
  The \bfem{space of (closed) $p$-forms of degree $n$} on an affine derived
  scheme $\Spec A$ is defined as
  \begin{align*}
\mathcal{A}^{p}_\K(A,n)
    &:= \mathrm{Map}_{\mathrm{dg}_k}(\K, (\wedge^p\LL_{A/\K})[n]), \\
\mathcal{A}^{p,cl}_\K(A,n)
    &:= \mathrm{Map}_{\epsilon-\mathrm{dg}_k}(\K(p)[-n-p],\DR(A/\K))
  \end{align*}
  in the $\infty$-category of mixed graded $\K$-modules \cite[Definition
  1.8]{PTVV}. For a derived Artin stack $F$, these are globalized via
  mapping spaces, see \cite[Definition 1.12]{PTVV}, such that
  \begin{equation*}
    \mathcal{A}^{p}_\K(F,n) :=
    \mathrm{Map}_{\dArt_\K}(F, \mathcal{A}^{p}_\K(-, n)), \quad \mathcal{A}^{p,cl}_\K(F,n) :=
    \mathrm{Map}_{\dArt_\K}(F, \mathcal{A}^{p,cl}_\K(-, n)).
  \end{equation*}  
\end{definition}
More specifically, $\omega \in \mathcal{A}^p(F,n)$ is a $d$-closed element of $DR(F)$ of weight $p$ and cohomological degree $n$. Similarly, $\omega \in \mathcal{A}^{p,cl}(F,n)$ is a $(d+d_{dR})$-closed element  of $DR(F)$ given as a formal series $\omega=\omega_p+\omega_{p+1} + \cdots$ such that $\omega_q$ has weight $q$ and degree $n-q+p$. For details, we refer to \cite{PTVV}.

Note also that the space $\mathcal{A}^{p,cl}(F,n)$  is usually more complicated than the space {$ \mathcal{A}^p(F,n)$  even if we make nice assumptions on $F$. Thanks to \cite[Prop.  1.14]{PTVV}, we have the following identification for the space $ \mathcal{A}^p(F,n)$, making it more tractable: 
\begin{equation*}
    \mathcal{A}^p(F,n)\simeq {\rm Map}_{L_{qcoh}(F)}(\mathcal{O}_F, \wedge^p\mathbb{L}_F[n]). 
\end{equation*}
 Thus, any 2-form of degree $n$  induces a morphism $\mathcal{O}_F \rightarrow  \wedge^2\mathbb{L}_F[n]$ in $D_{qcoh}(F)$, and hence, by duality, a morphism $\mathbb{T}_F \wedge \mathbb{T}_F \rightarrow \mathcal{O}_F[n]$. By adjunction, this gives the induced morphism $\Theta_{\omega} : \mathbb{T}_F \rightarrow \mathbb{L}_F [n]$, leading to the following definition:

\begin{definition}[Shifted Symplectic Structures]
  Following Pantev et al. \cite[Definition 1.18]{PTVV}, an
  \bfem{$n$-shifted symplectic structure} on a derived Artin stack
  $F$ is a point $\omega$ in the space $\mathbb{S}ymp(F, n)$. The
  space $\mathbb{S}ymp(F, n)$ is defined as the homotopy pullback
  \begin{equation*}
    \mathbb{S}ymp(F,n) := \mathcal{A}^{2,cl}(F,n)
    \times_{\mathcal{A}^2(F,n)}^h \mathcal{A}^2(F,n)^{nd}
  \end{equation*}
  where $\mathcal{A}^2(F,n)^{nd}$ is the subspace of non-degenerate
  2-forms, i.e., those whose underlying morphism $\Theta_{\omega}:
  \TT_F \to \LL_F[n]$ is an isomorphism in $D_{qcoh}(F)$.
\end{definition}

The shift $n$ compensates for the homological
degrees inherent in the tangent and cotangent complexes. This shift
distinguishes derived symplectic
geometry from its classical counterpart, admitting symplectic
structures on odd-dimensional moduli spaces \cite[Section 1.2]{PTVV}. 
Let us mention here two important existence results in the context of \cite{PTVV}.

\begin{theorem}[Calaque's Shifted Cotangent Stack Theorem] \label{thm:calaque}
  Let $X$ be a derived Artin stack. The $n$-shifted cotangent stack,
  defined via the relative derived spectrum as
  \begin{align*}
    T^*[n]X
    \simeq\mathbb{R}\mathbf{Spec}_X(\mathbb{L}\mathrm{Sym}_{\cO_X}(\TT_X[-n])),
  \end{align*}
  carries a canonical $n$-shifted symplectic structure \cite[Theorem
  2.2]{Calaque}.
\end{theorem}

Shifted cotangent stacks serve as local models for shifted symplectic
geometry, analogous to classical cotangent bundles \cite[Section 2.1]{Calaque}.

\begin{theorem}[PTVV Mapping Stack Theorem] \label{thm:ptvv}
  Let $F$ be a derived Artin stack equipped with an $n$-shifted
  symplectic structure, and let $Y$ be an $\cO$-compact derived stack
  equipped with an $\cO$-orientation of dimension $d$. Then the
  derived mapping stack  $\mathbb{R}\mathrm{Map}(Y, F)$ 
  has    a canonical $(n-d)$-shifted symplectic structure \cite[Theorem 2.5]{PTVV}.
\end{theorem} Here ``orientations," see \cite[Definition 2.4]{PTVV}, on derived stacks provide the homological equivalent of integration. This condition is a
prerequisite for pushing forward symplectic structures along mapping
spaces, analogous to integration along the fiber.
Furthermore, the mapping stack theorem generates examples of
shifted symplectic stacks. Moduli spaces of sheaves or
representations inherit shifted
symplectic structures from their domain curves or surfaces via the
mapping stack theorem. 
\vspace{0.1in}

Let us give some applications of the mapping stack theorem above.
\begin{example}
   Let $ G $ be an affine algebraic
    group  equipped with a nondegenerate invariant symmetric bilinear pairing  on its Lie algebra. Then the classifying stack is given as the functor   \[
\mathsf{B G} : Z \in \mathbf{dAff}_\K^{\mathrm{op}} \longmapsto 
\{\text{principal } G\text{-bundle on } Z\} 
\in \mathcal{S}.
\]Recall that $\mathsf{B G}$ has a \(2\)-shifted symplectic structure. Furthermore, there is an equivalence $\mathbb{T}_{\mathsf{B G}}\simeq \mathfrak{g}[1]$, where $\mathfrak{g}$ is the Lie algebra of G with the adjoint representation. 

Likewise, the \textit{moduli of perfect complexes} can be defined as the derived stack \[
\mathsf{Perf} : Z \in \mathbf{dAff}_\K^{\mathrm{op}} \longmapsto 
\{\text{perfect complexes on }   Z\} 
\in \mathcal{S}
\] such that it also carries a canonical \(2\)-shifted symplectic structure and that its tangent complex $\mathbb{T_{\mathsf{Perf}}}$ is equivalent to $E^{\vee} \otimes E[1]$, with $E$ the universal complex.

Combining these examples with the PTVV's mapping stack theorem above, let $Y$ be a smooth and projective Calabi-Yau $d$-fold, then by \cite{PTVV}:
\begin{enumerate}
  \item The derived Artin stack
    \( \mathrm{Perf}(Y) :=  \R\mathrm{Map}(Y, \mathrm{Perf}) \)
    of perfect complexes on $Y$ is \((2-d)\)-shifted symplectic.
    \item For a reductive group $G$, the derived Artin stack
    \( \mathrm{Bun}_G(Y) := \R\mathrm{Map}(Y, \mathsf{B G}) \)
    of principal $G$-bundles on $Y$ is \((2-d)\)-shifted symplectic.
\end{enumerate}
\end{example}

\begin{example}
    Fix the pair $(\mathsf{B G}, \omega)$ with $\omega \in \mathbb{S}ymp(\mathsf{B G}, 2)$. Given  a smooth and proper curve $C$, we call
    \begin{equation*}
        \mathrm{LocSys}_G(C) := \R\mathrm{Map} (C_{dR},\mathrm{BG})
    \end{equation*}
    the \bfem{moduli stack of flat $ G $-connections on $ C $}. Then Theorem \ref{thm:ptvv} also implies two important consequences:
    \begin{enumerate}
        \item $ \mathrm{LocSys}_G(C)$ is 0-shifted symplectic.
        \item There is a natural closed 2-form of degree 1 on $ \mathrm{Bun}_G(C) $, denoted by $\int_C ev^* \omega$, such that we have the identification \cite[Prop. 1.24]{Safronov2023} \begin{equation*}
        \mathrm{LocSys}_G(C)\simeq {\rm T}_{\int_C ev^* \omega}^*\mathrm{Bun}_G(C).
    \end{equation*}
            \end{enumerate} That is, $ \mathrm{LocSys}_G(C) $ can be equivalently seen as the \textit{twisted cotangent stack} of $ \mathrm{Bun}_G(C)$ with the twisting 2-form $ \int_C ev^* \omega \in \mathcal{A}^{2,cl}(\mathrm{Bun}_G(C), 1).$
\end{example}

Let us now list other key concepts, along with certain remarks and observations.
\begin{definition}[Isotropic and Lagrangian Structures]
  Let $(X, \omega)$ be an $n$-shifted symplectic derived stack. A map
  $f: L \to X$ is equipped with an \bfem{isotropic structure} if
  there is a specified path $h$ between $0$ and $f^*\omega$ in the
  Kan complex of
  closed 2-forms $\mathcal{A}^{2, cl}(L, n)$ \cite[Definition
  2.7]{PTVV}. 
  
  The Kan complex of isotropic structures is
  defined as $$Isot(f, \omega) := Path_{0,
  f^*\omega}(\mathcal{A}^{2,cl}(L, n)).$$ This
  null-homotopy induces a morphism $\Theta_h: \TT_L \to
  \LL_{L/X}[n-1]$. An isotropic structure is
  \bfem{Lagrangian} if $\Theta_h$ is an equivalence in the stable
  $\infty$-category $L_{qcoh}(L)$ \cite[Definition 2.8]{PTVV}.
\end{definition}

The requirement of a specified null-homotopy replaces the classical
condition of a vanishing differential form. In an $\infty$-category,
the space of null-homotopies carries higher categorical
structure; the choice of isotropic structure is part of the geometric
data, not a property.
\vspace{1in}

\begin{definition}[Shifted Contact Structures]
  Following Berktav \cite[Section 3.2]{Berktav1}, \cite[Definition 2.17]{Berktav2}, an \bfem{$n$-shifted contact structure} on a derived Artin stack $X$ consists of a line bundle $\cL$ on $X$, an $n$-shifted $1$-form $\alpha \in \pi_0 \mathbb{R}\Gamma(X, \LL_X \otimes^{\mathbb{L}}_{\cO_X} \cL[n])$, and a fiber sequence of perfect complexes built from the dual contraction map $\alpha^\vee \colon \TT_X \to \cL[n]$
  \begin{equation} \label{eq:contact_dist}
    \cK \to \TT_X \xrightarrow{\alpha^\vee} \cL[n]
  \end{equation}
  such that the underlying 2-form of the derived de Rham differential $d_{dR}\alpha$ induces an equivalence $\cK \xrightarrow{\sim} \cK^\vee \otimes^{\mathbb{L}}_{\cO_X} \cL[n]$ in $L_{qcoh}(X)$, rendering $\cK$ a symplectic complex.
\end{definition}

Shifted contact structures provide the
odd-dimensional counterpart to shifted symplectic
stacks. The line bundle accounts for the scaling ambiguity in contact
geometry, allowing the structure to be defined globally
across charts.

\begin{definition}[Derived Legendrian Structures]\label{definition:legendrian structures}
  Let $(X, \cL, \cK, \alpha)$ be an $n$-shifted contact derived stack. A map $f \colon L \to X$ is \bfem{isotropic} if it is equipped with a null-homotopy $f^*\alpha \simeq 0$. 

  An isotropic structure is called \bfem{Legendrian} if it induces a stable fiber sequence in $L_{qcoh}(L)$ of the form
  \begin{equation*}
    \cO_L \to \LL_{L/X} \otimes^{\mathbb{L}}_{\cO_L} f^*\cL[n-1] \to \TT_L.
  \end{equation*}
\end{definition}

\begin{remark}
  The \bfem{Legendrian} condition above is equivalent to the one given in \cite[Section 3.3]{Berktav3}. In \cite{Berktav3}, the isotropic lift $\TT_L \to f^*\cK$ induces a stable fiber sequence
  \begin{equation*}
    \TT_L \to f^*\cK \to \LL_L \otimes^{\mathbb{L}}_{\cO_L} f^*\cL[n].
  \end{equation*}
  Let $C := \mathrm{cofib}(\TT_L \to f^*\cK)$. The condition is $C \simeq \LL_L \otimes^{\mathbb{L}}_{\cO_L} f^*\cL[n]$.

  By definition of $\cK$, there is a sequence $f^*\cK \to f^*\TT_X \to f^*\cL[n]$. The relative tangent sequence for $f$ is $\TT_{L/X} \to \TT_L \to f^*\TT_X$.
  The pasting law for exact triangles of fibers (Lemma \ref{lem:pasting_law}) applied to $\TT_L \to f^*\cK \to f^*\TT_X$ produces a stable fiber sequence of cofibers
  \begin{equation*}
    C \to \TT_{L/X}[1] \to f^*\cL[n].
  \end{equation*}
  A shift by $-1$ and substitution of $C$ gives
  \begin{equation*}
    \LL_L \otimes^{\mathbb{L}}_{\cO_L} f^*\cL[n-1] \to \TT_{L/X} \to f^*\cL[n-1].
  \end{equation*}
  Taking the $\cO_L$-linear dual and tensoring with $f^*\cL[n-1]$ gives
  \begin{equation*}
    \cO_L \to \LL_{L/X} \otimes^{\mathbb{L}}_{\cO_L} f^*\cL[n-1] \to \TT_L.
  \end{equation*}
\end{remark}

The Legendrian condition establishes the derived geometric analogue
for maximally isotropic submanifolds. As in the classical
theory, derived Legendrian structures are governed by the
underlying contact distribution, providing the geometric foundation
for studying boundary conditions in topological field theories.
\section{Transversality}
\subsection{Classical Transversality} \label{sec:classical}

\begin{definition}[Liouville Vector Field]
  Let $(W, \omega)$ be a symplectic manifold. A vector field $Y$ on
  $W$ is called a \bfem{Liouville vector field} if its Lie
  derivative satisfies $\mathcal{L}_Y \omega = \omega$.
\end{definition}

\begin{lemma}[{\cite[Lemma 1.4.5]{Geiges}}] \label{lem:classical_geiges}
  Let $(W, \omega)$ be a symplectic manifold of dimension $2n$, and
  let $Y$ be a Liouville vector field on $W$. If $M \subset W$ is a
  smooth hypersurface transverse to $Y$, then the restriction $\alpha
  := (\iota_Y \omega)|_M$ is a contact form on $M$.
\end{lemma}
\begin{proof}
  By Cartan's magic formula and the assumption that the symplectic
  form is closed ($d\omega = 0$), the following holds
  \begin{equation*}
    d(\iota_Y \omega) = \mathcal{L}_Y \omega - \iota_Y d\omega = \omega.
  \end{equation*}
  Therefore, $d\alpha = \omega|_M$. Showing that $\alpha$ is a
  contact form requires verifying that $\alpha \wedge (d\alpha)^{n-1}$
  is a nowhere vanishing volume form on $M$.

  On the symplectic manifold $W$, the top-dimensional form $\omega^n$
  is a volume form. Contracting this with the Liouville vector field $Y$ we obtain
  \begin{equation*}
    \iota_Y(\omega^n) = n (\iota_Y \omega) \wedge \omega^{n-1}.
  \end{equation*}
  Because $Y$ is transverse to the hypersurface $M$, the tangent
  space admits the pointwise coordinate splitting $T_x W \cong T_x M
  \oplus \mathbb{R}Y_x$.
  
  Evaluating the top-form $\omega^n$ on the frame $(Y_x, e_1, \dots,
  e_{2n-1})$, where the vectors $e_i$ span the local tangent space
  $T_x M$, shows that the transverse contraction
  $\iota_Y(\omega^n)|_M$ defines a nowhere vanishing volume form on
  $M$. Restricting the identity to $M$ preserves this
  property, yielding
  \begin{equation*}
    n \alpha \wedge (d\alpha)^{n-1} \neq 0,
  \end{equation*}
  proving $\alpha$
  is a contact form on $M$.
\end{proof}

The classical proof relies on point-set transversality
(the splitting of the tangent space) and the pointwise evaluation
of differential forms on topological frames. In derived moduli
theory, intersections are rarely transverse and tangent
spaces jump in rank, requiring a homological replacement for
this coordinate argument.
\subsection{The Derived Setup and Statement} \label{sec:setup}

The classical Liouville vector field is replaced by a group action in
the derived setting.

\begin{definition}[Lie Algebra of a Group Stack]
  For an affine smooth group scheme $G$ with Lie algebra
  $\mathfrak{g}$, the cotangent complex of the classifying stack $BG$
  is $\mathfrak{g}^\vee[-1]$. A $G$-action on a stack $\widetilde{X}$
  induces an infinitesimal Lie algebra action morphism
  $\mathfrak{g} \otimes_\K \cO_{\widetilde{X}} \to
  \TT_{\widetilde{X}}$ in $L_{qcoh}(\widetilde{X})$ \cite[Section 1.2]{PTVV}.
\end{definition}

\begin{lemma}[Derived Atiyah Sequence] \label{lem:atiyah}
  Let $G$ be a smooth affine group scheme with Lie algebra $\mathfrak{g}$. For a principal $G$-torsor $p \colon \widetilde{X} \to X$ in $\dArt_\K$, the $G$-action induces a stable fiber sequence in $L_{qcoh}(\widetilde{X})$
  \begin{equation*}
    \mathbb{L}p^*\LL_{X} \to \LL_{\widetilde{X}} \to \mathfrak{g}^\vee \otimes_\K \cO_{\widetilde{X}}.
  \end{equation*}
\end{lemma}
\begin{proof}
  By \cite[Definition 1.4.1.15]{HAG-II}, the morphism $p \colon \widetilde{X} \to X$ has a relative cotangent complex. By \cite[Lemma 1.4.1.16(1)]{HAG-II}, for any affine point $x \colon \Spec(A) \to \widetilde{X}$, there is a homotopy cofiber sequence of stable $A$-modules
  \begin{equation} \label{eq:local_transitivity}
    \LL_{X, p \circ x} \to \LL_{\widetilde{X}, x} \to \LL_{\widetilde{X}/X, x}.
  \end{equation}
  
  By \cite[Definitions 1.4.1.7 and 1.4.1.15]{HAG-II}, the pointwise complexes are the pullbacks of the global cotangent complexes, implying $$\LL_{\widetilde{X}, x} \simeq x^* \LL_{\widetilde{X}}, \quad \LL_{\widetilde{X}/X, x} \simeq x^* \LL_{\widetilde{X}/X}, \; \text{ and } \ \LL_{X, p \circ x} \simeq x^* \mathbb{L}p^* \LL_X.$$ 
  
  Substituting these identifications into \eqref{eq:local_transitivity} gives a stable fiber sequence in $L_{qcoh}(\Spec(A))$ for every affine point $x$
  \begin{equation*}
    x^* \mathbb{L}p^* \LL_X \to x^* \LL_{\widetilde{X}} \to x^* \LL_{\widetilde{X}/X}.
  \end{equation*}
  
  The stable $\infty$-category $L_{qcoh}(\widetilde{X})$ is equivalent to the homotopy limit of $L_{qcoh}(\Spec(A))$ over all affine points mapping to $\widetilde{X}$. Pointwise exactness establishes the global transitivity fiber sequence in $L_{qcoh}(\widetilde{X})$
  \begin{equation} \label{eq:standard_transitivity}
    \mathbb{L}p^*\LL_{X} \to \LL_{\widetilde{X}} \to \LL_{\widetilde{X}/X}.
  \end{equation}
  
  Because $p$ is a principal $G$-torsor, the action map $\rho \colon G \times \widetilde{X} \to \widetilde{X}$ and the projection $\mathrm{pr}_2 \colon G \times \widetilde{X} \to \widetilde{X}$ define a homotopy pullback square in $\dArt_\K$
  \[
    \begin{tikzcd}
      G \times \widetilde{X} \arrow[r, "\rho"] \arrow[d, "\mathrm{pr}_2"'] & \widetilde{X} \arrow[d, "p"] \\
      \widetilde{X} \arrow[r, "p"'] & X
    \end{tikzcd}
  \]
  
  By \cite[Lemma 1.4.1.16(2)]{HAG-II}, base change along $p$ gives the equivalence
  \begin{equation*}
    \mathbb{L}\rho^* \LL_{\widetilde{X}/X} \simeq \LL_{G \times \widetilde{X} / \widetilde{X}}.
  \end{equation*}
  
  The morphism $\mathrm{pr}_2$ is the base change of the structural morphism $G \to \Spec(\K)$ along $\pi_{\widetilde{X}} \colon \widetilde{X} \to \Spec(\K)$. Applying \cite[Lemma 1.4.1.16(2)]{HAG-II} gives the equivalence $\LL_{G \times \widetilde{X} / \widetilde{X}} \simeq \mathbb{L}\mathrm{pr}_1^* \LL_{G}$, yielding
  \begin{equation*}
    \mathbb{L}\rho^* \LL_{\widetilde{X}/X} \simeq \mathbb{L}\mathrm{pr}_1^* \LL_{G}.
  \end{equation*}
  
  Pulling back this equivalence along the section $s = (e, \mathrm{id}_{\widetilde{X}}) \colon \widetilde{X} \to G \times \widetilde{X}$ gives
  \begin{equation*}
    \mathbb{L}s^* \mathbb{L}\rho^* \LL_{\widetilde{X}/X} \simeq \mathbb{L}s^* \mathbb{L}\mathrm{pr}_1^* \LL_{G}.
  \end{equation*}
  
  The section $s$ satisfies the composition identities $\rho \circ s = \mathrm{id}_{\widetilde{X}}$ and $\mathrm{pr}_1 \circ s = e \circ \pi_{\widetilde{X}}$, shown in the commutative diagrams
  \[
    \begin{tikzcd}
        \widetilde{X} \arrow[r, "s"] \arrow[dr, "\mathrm{id}_{\widetilde{X}}"'] & G \times \widetilde{X} \arrow[d, "\rho"] \\
         & \widetilde{X}
    \end{tikzcd}
    \quad \quad
    \begin{tikzcd}
        \widetilde{X} \arrow[r, "s"] \arrow[d, "\pi_{\widetilde{X}}"'] & G \times \widetilde{X} \arrow[d, "\mathrm{pr}_1"] \\
        \Spec(\K) \arrow[r, "e"'] & G
    \end{tikzcd}
  \]
  
  These identities reduce the pullback equivalence to
  \begin{equation*}
    \LL_{\widetilde{X}/X} \simeq \mathbb{L}\pi_{\widetilde{X}}^* \mathbb{L}e^* \LL_{G}.
  \end{equation*}
  
  For a smooth affine group scheme $G$, the pullback of the cotangent complex along the identity section gives $\mathbb{L}e^* \LL_{G} \simeq \mathfrak{g}^\vee$. This establishes the equivalence $\LL_{\widetilde{X}/X} \simeq \mathfrak{g}^\vee \otimes_{\mathbb{K}} \cO_{\widetilde{X}}$. Substituting this identification into \eqref{eq:standard_transitivity} produces the stable fiber sequence $$\mathbb{L}p^*\LL_{X} \to \LL_{\widetilde{X}} \to \mathfrak{g}^\vee \otimes_{\mathbb{K}} \cO_{\widetilde{X}}.$$
    \end{proof}

\begin{corollary} \label{cor:fundamental_morphism}
  Let $p \colon \widetilde{X} \to X$ be a principal $\Gm$-torsor in $\dArt_\K$. The standard left-invariant vector field on $\Gm$ determines a fundamental morphism $Y \colon \cO_{\widetilde{X}} \to \TT_{\widetilde{X}}$ in $L_{qcoh}(\widetilde{X})$ such that the relative cotangent sequence takes the form
  \begin{equation*}
    \LL p^* \LL_X \xrightarrow{(\TT p)^\vee} \LL_{\widetilde{X}} \xrightarrow{Y^\vee} \cO_{\widetilde{X}}.
  \end{equation*}
\end{corollary}
\begin{proof}
  The proof of Lemma \ref{lem:atiyah} establishes the equivalence $\LL_{\widetilde{X}/X} \simeq \mathfrak{g}^\vee \otimes_\K \cO_{\widetilde{X}}$. Taking the derived $\cO_{\widetilde{X}}$-linear dual of this equivalence gives an isomorphism
  \begin{equation*}
    \TT_{\widetilde{X}/X} \simeq \mathfrak{g} \otimes_\K \cO_{\widetilde{X}}.
  \end{equation*}
  The relative tangent complex sits in the transitivity fiber sequence $\TT_{\widetilde{X}/X} \to \TT_{\widetilde{X}} \xrightarrow{\TT p} \LL p^* \TT_X$. Composing the isomorphism with the canonical morphism $\TT_{\widetilde{X}/X} \to \TT_{\widetilde{X}}$ defines the infinitesimal action morphism
  \begin{equation*}
    a \colon \mathfrak{g} \otimes_\K \cO_{\widetilde{X}} \to \TT_{\widetilde{X}}.
  \end{equation*}
  The multiplicative group scheme $\Gm$ has the coordinate algebra $\K[t, t^{-1}]$. The tangent space at the identity is the $1$-dimensional $\K$-vector space generated by the derivation $\partial = t \partial_t$. This generator defines an isomorphism of $\K$-vector spaces $\K \xrightarrow{\sim} \mathfrak{g}$. The fundamental morphism $Y$ is defined as the composition
  \begin{equation*}
    Y \colon \cO_{\widetilde{X}} \simeq \K \otimes_\K \cO_{\widetilde{X}} \xrightarrow{\sim} \mathfrak{g} \otimes_\K \cO_{\widetilde{X}} \xrightarrow{a} \TT_{\widetilde{X}}.
  \end{equation*}
  Taking the derived dual of this composition we obtain
  \begin{equation*}
    Y^\vee \colon \LL_{\widetilde{X}} \xrightarrow{a^\vee} \mathfrak{g}^\vee \otimes_\K \cO_{\widetilde{X}} \xrightarrow{\sim} \cO_{\widetilde{X}}.
  \end{equation*}
  The morphism $a^\vee$ is the canonical projection $\LL_{\widetilde{X}} \to \LL_{\widetilde{X}/X}$ in the cotangent transitivity sequence, post-composed with the equivalence from Lemma \ref{lem:atiyah}. Substituting this identification into the standard transitivity sequence gives the stable fiber sequence
  \begin{equation*}
    \LL p^* \LL_X \xrightarrow{(\TT p)^\vee} \LL_{\widetilde{X}} \xrightarrow{Y^\vee} \cO_{\widetilde{X}}.
  \end{equation*}
\end{proof}

The $\Gm$-action allows the formulation of the contact space
without local analytic arguments \cite[Section 4.2]{Berktav1}.

\begin{definition}[Derived Liouville Condition]
  The $\Gm$-action has \textbf{weight 1} with respect to
  $\omega_{\widetilde{X}}$ if $\omega_{\widetilde{X}}$ is a weight 1 eigenform.
  Equivalently, the derived pullback of the action map
  $\rho: \Gm \times \widetilde{X} \to \widetilde{X}$ satisfies $\rho^*
  \omega_{\widetilde{X}} \simeq z \cdot \pi_{\widetilde{X}}^*
  \omega_{\widetilde{X}}$,
  where $z$ is the character coordinate on $\Gm$
  \cite[Section 2.1.1]{Calaque}.
\end{definition}
 

In classical geometry, the Liouville vector field dilates the
symplectic form. The weight 1 condition is the algebraic
translation of this dilation, rendering the symplectic form
homogeneous of degree 1. This homogeneity ensures the contact form
survives the quotient process.
The derived analogue of the classical transversality lemma can then be stated as follows.

\begin{theorem}[Derived Analogue of Classical Transversality] \label{thm:derived transversality}
  Let $(\widetilde{X}, \omega_{\widetilde{X}})$ be an $n$-shifted symplectic
  derived Artin stack equipped with a $\Gm$-action of weight 1. Let
  $X$ be the stack quotient $X := [\widetilde{X} / \Gm]$, and let $p \colon
  \widetilde{X} \to X$ be the projection. Then $X$ inherits an
  $n$-shifted contact structure.
\end{theorem}

\begin{remark} \label{rmk:darboux_atlas}
  For $n < 0$, the descended contact structure
  on $X$ locally admits a contact Darboux atlas \cite[Theorem 3.7]{Berktav2}.
\end{remark}

Taking the quotient of a homogenous symplectic space descends the
symplectic data to a contact structure, avoiding a transverse
hypersurface \cite[Section 2.26]{Berktav2}.

\subsection{Proof of Theorem \ref{thm:A} and Analysis of Transversality}\label{sec:proof}

\begin{observation}
  \label{prop:transversality_failure}
  Transversality between a global vector field and a closed immersion of virtual codimension $1$ fails for proper derived Artin stacks, provided that the zero locus of the vector field intersects the closed immersion.
\end{observation}
\begin{proof}
  Let $\widetilde{X}$ be a proper derived Artin stack. Let $Y \colon \cO_{\widetilde{X}} \to \TT_{\widetilde{X}}$ be a global vector field in $L_{qcoh}(\widetilde{X})$. Let $i \colon H \to \widetilde{X}$ be a closed immersion. The derived normal complex $\mathcal{N}_{H/\widetilde{X}}$ is the cofiber in the tangent transitivity sequence
  \begin{align*}
    \TT_H \xrightarrow{di} i^*\TT_{\widetilde{X}} \to \mathcal{N}_{H/\widetilde{X}}.
  \end{align*}
  Classical transversality requires the composite morphism $c \colon \cO_H \xrightarrow{i^*Y} i^*\TT_{\widetilde{X}} \to \mathcal{N}_{H/\widetilde{X}}$ to be an equivalence in $L_{qcoh}(H)$. This equivalence establishes a homological splitting of the tangent sequence, generating a rank-$1$ free summand on $H$.

  The derived zero locus $Z(Y)$ is the homotopy pullback of the section $Y$ and the zero section $0 \colon \cO_{\widetilde{X}} \to \TT_{\widetilde{X}}$ in the $\infty$-category of derived stacks. If $\widetilde{X}$ possesses a non-vanishing virtual Euler class, the Poincar\'{e}-Hopf index theorem ensures $Z(Y)$ is non-empty.

  Suppose there exists an $R$-point $x \colon \mathrm{Spec}(R) \to H \times_{\widetilde{X}}^h Z(Y)$ in the derived intersection of the hypersurface and the zero locus. By the definition of the homotopy pullback $Z(Y)$, the restricted vector field $x^*Y \colon \cO_{\mathrm{Spec}(R)} \to x^*\TT_{\widetilde{X}}$ is equipped with a null-homotopy $x^*Y \simeq 0$.

  This null-homotopy trivializes the pulled-back composite morphism $x^*c$. Since $c$ is an equivalence, the pullback $x^*c$ must be an equivalence in $L_{qcoh}(\mathrm{Spec}(R))$. An equivalence between rank-$1$ free modules cannot be null-homotopic over a non-trivial ring $R$. The existence of zeros on $H$ obstructs the equivalence $c$. Global transversality fails.  
\end{proof}

\begin{remark}
  \label{rmk:slice_vs_quotient}
  The stack quotient $X \simeq |B_\bullet(\widetilde{X}, \mathbb{G}_m)|$ circumvents the transversality failure by absorbing the zeros of the vector field into its stabilizer structure. The vector field $Y$ is the infinitesimal generator of the $\mathbb{G}_m$-action on $\widetilde{X}$. The derived zero locus $Z(Y)$ is equivalent to the derived fixed point stack $\widetilde{X}^{\mathbb{G}_m}$. Instead of producing a topological singularity in the quotient, the fixed points acquire $\mathbb{G}_m$ as their stabilizer group in the resulting $\infty$-groupoid.

The projection $p \colon \widetilde{X} \to X$ defines the tangent fiber sequence $\cO_{\widetilde{X}} \xrightarrow{Y} \TT_{\widetilde{X}} \to \mathbb{L}p^* \TT_X$ in $L_{qcoh}(\widetilde{X})$. For an $R$-point $x \colon \mathrm{Spec}(R) \to Z(Y)$, the vanishing of the vector field corresponds to the null-homotopy $x^*Y \simeq 0$. This null-homotopy induces the splitting $$x^*\mathbb{L}p^*\TT_X \simeq x^*\TT_{\widetilde{X}} \oplus \cO_{\mathrm{Spec}(R)}[1] \quad \text{ in } L_{qcoh}(\mathrm{Spec}(R)).$$ The homological failure of transversality is converted into the shifted Lie algebra $\cO_{\mathrm{Spec}(R)}[1]$ of the stabilizer group. The complex $\mathbb{L}p^*\TT_X$ remains perfect globally because $\mathrm{Perf}(\widetilde{X})$ is a stable subcategory of $L_{qcoh}(\widetilde{X})$. This mechanism preserves geometric dualizability without imposing a global homological splitting on the tangent sequence.
\end{remark}

\begin{proof}[\textbf{Proof of Theorem \ref{thm:derived transversality}}]
  The proof proceeds in four steps. Let us outline each step: 
  
  We first trivialize the relative cotangent complex of the principal $\Gm$-torsor $p \colon \widetilde{X} \to X$ using the infinitesimal $\Gm$-action. We then construct an $n$-shifted $1$-form $\alpha_{\widetilde{X}}$ on $\widetilde{X}$ by contracting the symplectic form $\omega_{\widetilde{X}}$ with the fundamental vector field of the action. We prove this $1$-form is horizontal with respect to $p$. The horizontality and the weight 1 equivariant structure of $\omega_{\widetilde{X}}$ allow the descent of $\alpha_{\widetilde{X}}$ to a twisted $1$-form $\alpha$ on the quotient stack $X$. Finally, we define the derived contact distribution as the homotopy fiber of the descended $1$-form and deduce its non-degeneracy from the underlying symplectic structure on $\widetilde{X}$.
\vspace{5pt}

  \textbf{\textit{Step 1: The Relative Cotangent Sequence.}}
  Let $X$ denote the stack quotient $[\widetilde{X} / \Gm]$. By definition, the projection $p \colon \widetilde{X} \to X$ is a principal $\Gm$-torsor. By Corollary \ref{cor:fundamental_morphism}, the torsor determines a fundamental morphism $Y$ and induces the relative cotangent sequence
  \begin{equation*}
    \LL p^* \LL_X \xrightarrow{(\TT p)^\vee} \LL_{\widetilde{X}} \xrightarrow{Y^\vee} \cO_{\widetilde{X}}.
  \end{equation*}

  \textbf{\textit{Step 2: Construction of the 1-Form.}}
  The derived Artin stack $\widetilde{X}$ is equipped with an
  $n$-shifted symplectic structure $\omega_{\widetilde{X}}$. By \cite[Definition 1.18]{PTVV}, the
  non-degeneracy of $\omega_{\widetilde{X}}$ gives an equivalence of perfect
  complexes $\Theta_{\omega} \colon \TT_{\widetilde{X}} \xrightarrow{\sim}
  \LL_{\widetilde{X}}[n]$.
  Post-composing the fundamental morphism $Y \colon \cO_{\widetilde{X}} \to
  \TT_{\widetilde{X}}$ with this equivalence yields a morphism in $L_{qcoh}(\widetilde{X})$,
  identified via the Yoneda lemma
  $\mathbb{R}\Gamma(\widetilde{X}, E) \simeq
  \mathrm{Map}_{L_{qcoh}(\widetilde{X})}(\cO_{\widetilde{X}}, E)$,
  with an $n$-shifted $1$-form global section
  $\alpha_{\widetilde{X}} \in \pi_0 \mathbb{R}\Gamma(\widetilde{X},
  \LL_{\widetilde{X}}[n])$ \cite[Section 4.2]{Berktav1} such that
  \begin{equation*}
    \alpha_{\widetilde{X}} := \Theta_{\omega} \circ Y.
  \end{equation*}
  By the defining equivalence, this global section corresponds
  to a contraction morphism
  $\alpha_{\widetilde{X}}^\vee \colon \TT_{\widetilde{X}} \to \cO_{\widetilde{X}}[n]$.

  The composition $\cO_{\widetilde{X}} \xrightarrow{Y} \TT_{\widetilde{X}} \xrightarrow{\alpha_{\widetilde{X}}^\vee} \cO_{\widetilde{X}}[n]$ corresponds to the derived self-pairing of $Y$ with respect to $\omega_{\widetilde{X}}$. In the graded mixed derived de Rham complex \cite[Section 1.4.1]{CPTVV}, the graded commutativity of the contraction operator implies that the interior product $\iota_Y \iota_Y \omega_{\widetilde{X}}$ admits a canonical null-homotopy. This defines an equivalence $\omega_{\widetilde{X}}(Y, Y) \simeq 0$ within the mapping space $\mathrm{Map}_{L_{qcoh}(\widetilde{X})}(\cO_{\widetilde{X}}, \cO_{\widetilde{X}}[n])$.

  Applying the contravariant derived internal Hom functor
  $\mathbb{R}\underline{\mathcal{H}om}( - , \cO_{\widetilde{X}}[n])$
  to the exact triangle of tangent complexes $\cO_{\widetilde{X}} \xrightarrow{Y}
  \TT_{\widetilde{X}} \xrightarrow{\mathbb{T}p} \mathbb{L}p^*\TT_X \xrightarrow{+1}$
  gives the dual exact triangle 
  \begin{equation*}
    \mathbb{L}p^*\LL_X[n] \xrightarrow{(\mathbb{T}p)^\vee}
    \LL_{\widetilde{X}}[n] \xrightarrow{Y^\vee} \cO_{\widetilde{X}}[n] \xrightarrow{+1}
  \end{equation*}
  in $L_{qcoh}(\widetilde{X})$.
  Evaluating global sections by applying the functor $\mathrm{Map}_{L_{qcoh}(\widetilde{X})}(\cO_{\widetilde{X}}, -)$ to this exact triangle generates a long exact sequence in cohomology containing the segment
  \begin{equation*}
    \cdots \to \pi_0 \mathrm{Map}(\cO_{\widetilde{X}}, \mathbb{L}p^*\LL_X[n]) \xrightarrow{(\mathbb{T}p)^\vee_*} \pi_0 \mathrm{Map}(\cO_{\widetilde{X}}, \LL_{\widetilde{X}}[n]) \xrightarrow{Y^\vee_*} \pi_0 \mathrm{Map}(\cO_{\widetilde{X}}, \cO_{\widetilde{X}}[n]) \to \cdots
  \end{equation*}
  The 1-form $\alpha_{\widetilde{X}}$ defines a class in $\pi_0 \mathrm{Map}(\cO_{\widetilde{X}}, \LL_{\widetilde{X}}[n])$. Its image under $Y^\vee_*$ computes the self-pairing $\omega_{\widetilde{X}}(Y, Y)$. The specified null-homotopy $\omega_{\widetilde{X}}(Y, Y) \simeq 0$ forces $Y^\vee_*(\alpha_{\widetilde{X}}) = 0$. By exactness of the sequence, the kernel of $Y^\vee_*$ equals the image of $(\mathbb{T}p)^\vee_*$. This guarantees the existence of a lift of $\alpha_{\widetilde{X}}$ to a section in $\pi_0 \mathrm{Map}(\cO_{\widetilde{X}}, \mathbb{L}p^* \LL_X[n])$.
\vspace{5pt}

  \textbf{\textit{Step 3: Equivariant Descent.}}
  To establish that the 1-form $\alpha_{\widetilde{X}}$ descends to the quotient stack $X$, we must determine its equivariant structure. The stable $\infty$-category of perfect complexes on the geometric realization $X \simeq \varinjlim_{\Delta^{\mathrm{op}}} B_\bullet(\Gm, \widetilde{X})$ is equivalent to the stable $\infty$-category of $\Gm$-equivariant perfect complexes on $\widetilde{X}$. The $\Gm$-action lifts to the cotangent complex $\LL_{\widetilde{X}}$. 
  By the premise of the theorem, the symplectic form $\omega_{\widetilde{X}}$ is an eigenform of weight 1 with respect to this action. Since $\Gm$ is abelian, the fundamental morphism $Y$ inducing the action is invariant under the adjoint action, as it originates from the Lie algebra of $\Gm$, giving it a weight 0 equivariant structure. Since the equivalence $\Theta_\omega$ intertwines the equivariant structures up to a weight 1 twist, post-composing the weight 0 morphism $Y$ with $\Theta_\omega$ gives a 1-form $\alpha_{\widetilde{X}}$ with a weight 1 equivariant structure \cite[Section 1.2]{PTVV}.

  Because $\alpha_{\widetilde{X}}$ is both horizontal, as established in \textbf{\textit{Step 2}}, and invariant up to a weight 1 twist, it descends to the quotient stack. Let $\cL$ be the line bundle on $X$ associated to the weight 1 representation of $\Gm$. By the derived descent theory of $\Gm$-torsors, this $\Gm$-equivariant horizontal global section descends to a global section on $X$ taking values in this twisted representation. This yields the descended 1-form
  \begin{equation*}
    \alpha \in \pi_0 \mathbb{R}\Gamma(X, \LL_X \otimes^{\mathbb{L}}_{\cO_X} \cL[n]).
  \end{equation*}

\textbf{\textit{Step 4: Constructing the Contact Complex and Proving Non-Degeneracy.}}
  On the quotient stack $X$, the equivalence
  $\LL_X \otimes^{\mathbb{L}}_{\cO_X} \cL[n] \simeq
  \mathbb{R}\underline{\mathcal{H}om}(\TT_X, \cL[n])$ identifies
  the descended 1-form $\alpha$ with a contraction morphism
  $\alpha^\vee \colon \TT_X \to \cL[n]$
  within $L_{qcoh}(X)$. The derived contact distribution
  $\cK \in L_{qcoh}(X)$ is defined as the homotopy fiber of this map, yielding
  the stable fiber sequence 
  \begin{equation} \label{eq:contact_dist}
    \cK \to \TT_X \xrightarrow{\alpha^\vee} \cL[n]
  \end{equation}
  in $L_{qcoh}(X)$ \cite[Definition 3.5]{Berktav1}. 

  By pulling the fiber sequence \eqref{eq:contact_dist} back to
  $\widetilde{X}$ via $\mathbb{L}p^*$,
  the twisted line bundle $\mathbb{L}p^*\cL[n]$ trivializes to
  $\cO_{\widetilde{X}}[n]$.
  The derived pullback of $\cK$, denoted $\mathbb{L}p^*\cK$, fits
  into the fiber sequence
  \begin{equation*}
    \mathbb{L}p^*\cK \to \mathbb{L}p^*\TT_X
    \xrightarrow{\mathbb{L}p^*\alpha^\vee} \cO_{\widetilde{X}}[n].
  \end{equation*}

  The underlying symplectic form provides the equivalence
  $\Theta_\omega \colon \TT_{\widetilde{X}} \xrightarrow{\sim}
  \LL_{\widetilde{X}}[n]$. The orthogonal complement $Y^\perp$ is the
  homotopy fiber of the contraction map
  $\alpha_{\widetilde{X}}^\vee \colon \TT_{\widetilde{X}} \to
  \cO_{\widetilde{X}}[n]$. As established
  in \textbf{\textit{Step 2}}, the pairing evaluated on the fundamental morphism yields
  the path $\omega_{\widetilde{X}}(Y, Y) \simeq 0$. This path defines the
  null-homotopy $\alpha_{\widetilde{X}}^\vee \circ Y \simeq 0$. This datum lifts the morphism
  $Y$ to the homotopy fiber $Y^\perp$, showing $\cO_{\widetilde{X}}$
  as an isotropic object within the symplectic $\infty$-category and
  yielding the lift $\cO_{\widetilde{X}} \to Y^\perp$.

  The morphism $\alpha_{\widetilde{X}}^\vee$ factors as the composition $\mathbb{L}p^*\alpha^\vee \circ \mathbb{T}p$. 
  The homotopy fiber of $\TT p\colon \TT_{\widetilde{X}}\to \LL p^*\TT_{X}$ is the relative tangent complex $\TT_{\widetilde{X}/X}$ which is by \textbf{\textit{Step 1}} equivalent to $\cO_{\widetilde{X}}$. The homotopy fiber of $\alpha_{\widetilde{X}}^\vee$ is $Y^\perp$ by definition and the homotopy fiber of $\LL p^*\alpha^\vee\colon \LL p^* \TT_X \to \cO_{\widetilde{X}}[n]$ is the pullback of the contact distribution $\LL p^*\cK$. Applying Lemma \ref{lem:pasting_law} to $\LL p^* \alpha^\vee \circ \TT p$ yields the fiber sequence
  \begin{equation*}
    \cO_{\widetilde{X}}\to Y^\perp \to \LL p^*\cK   
  \end{equation*}
  making $\LL p^*\cK$ the homotopy cofiber of the isotropic lift $\cO_{\widetilde{X}}\to Y^\perp$.
  
  In the stable $\infty$-category of perfect complexes, the homotopy cofiber of an isotropic object lifted to
  its orthogonal complement possesses a non-degenerate pairing \cite[Section 1.1.2]{LurieHA}.
  The derived symplectic reduction of the isotropic morphism $Y$ guarantees the equivalence $\Theta_\omega$ descends to an equivalence
  $(\mathbb{L}p^*\cK) \xrightarrow{\sim} (\mathbb{L}p^*\cK)^\vee[n]$
  on the total space. 

  Within the mixed graded derived de Rham complex \cite[Section 1.4.1]{CPTVV}, the derived de Rham differential is natively encoded by the mixed structure $\epsilon$. Because the symplectic form $\omega_{\widetilde{X}}$ is an eigenform of weight 1, the $\Gm$-equivariant descent ensures that the mixed differential of the descended 1-form $\alpha$ recovers the symplectic form upon pullback. Therefore, the underlying 2-form of the derived de Rham differential $p^*(d_{dR}\alpha)$ is identified with $\omega_{\widetilde{X}}$.
  
  This equivariant non-degeneracy guarantees that the derived de Rham differential $d_{dR}\alpha$ induces the equivalence $\cK \xrightarrow{\sim} \cK^\vee \otimes^{\mathbb{L}}_{\cO_X} \cL[n]$ on the quotient $X$ after applying descent.
  
We then complete the proof of Theorem \ref{thm:derived transversality}, and hence that of Theorem \ref{thm:A}.
\end{proof}

\section{Derived Legendrian Intersection Theorem}\label{sec:legendrian}
Before establishing the intersection theorem, the categorical
relationship between Legendrians and the derived symplectification
functor must be formalized. 

Let us recall our setting and fix our notation: the \bfem{${n}$-shifted contact structure} on a
derived Artin stack $X$ is defined by a contact line bundle $\cL$, an
$n$-shifted $1$-form $\alpha$, and a non-degenerate contact
distribution \[\cK \simeq \mathrm{fib}(\TT_X \xrightarrow{\alpha^\vee}
\cL[n]) \ \text{ in } L_{qcoh}(X).\]

Following \cite[Definition 4.3 and Proposition 4.4]{Berktav1}, the \bfem{derived symplectification} of $X$, denoted by $\widetilde{X}$, is defined as the total space of the $\Gm$-bundle associated with the contact line bundle $\cL$. 
The projection $p \colon \widetilde{X} \to X$ classifies trivializations of $\cL$, forming a principal $\Gm$-bundle in $\dArt_\K$ \cite[Section 1.3.5]{HAG-II}. 
By \cite[Theorem 4.7]{Berktav1} for derived schemes and \cite[Theorem 3.9]{Berktav2} for derived Artin stacks, $\widetilde{X}$ is equipped with an $n$-shifted symplectic structure $\omega_{\widetilde{X}}$. 
This symplectic form is an eigenform of weight 1 with respect to the $\Gm$-action.

A \bfem{Legendrian morphism} $f \colon L \to X$ provides a specified
null-homotopy $f^*\alpha \simeq 0$ and induces a stable fiber sequence
\begin{equation*}
  \cO_L \to \LL_{L/X} \otimes^{\mathbb{L}}_{\cO_L} f^*\cL[n-1] \to \TT_L
\end{equation*}
of perfect complexes in $L_{qcoh}(L)$ (cf. Definition \ref{definition:legendrian structures}).

The key is that derived contact geometry links these
Legendrian structures to Lagrangian structures via $\Gm$-equivariant
lifts along the symplectification projection $p$.
This lifting property translates intersection problems in contact
moduli to equivariant intersection problems in symplectic moduli, and hence establishes the shifted contact structure on the derived
intersection of two Legendrian stacks. More specifically, we have: 

\begin{theorem}  \label{thm:Leg intersection}
  Let $X$ be an $n$-shifted contact derived Artin stack. Let $L_1$ and
  $L_2$ be two derived Artin stacks equipped with Legendrian morphisms
  \[\begin{tikzcd}
                    & L_2 \arrow[d, "f_2"] \\
L_1 \arrow[r, "f_1"] & X.                   
\end{tikzcd}\]
  Then the derived
  fiber product $Z := L_1 \times_X^h L_2$ carries an $(n-1)$-shifted
  contact structure.
\end{theorem}
\vspace{0.1in}

\begin{proof}
  The proof proceeds in four steps utilizing the derived symplectification functor.

  \vspace{5pt}

  \textbf{\textit{Step 1: Equivariant lifting of the intersecting stacks.}}
  Let $\widetilde{X}$ be the derived symplectification of $X$. The domain stacks $\widetilde{L}_i := L_i \times_X^h \widetilde{X}$ are principal $\Gm$-bundles over $L_i$ via the projections $q_i \colon \widetilde{L}_i \to L_i$. The derived fiber product induces $\Gm$-equivariant morphisms $\widetilde{f}_i \colon \widetilde{L}_i \to \widetilde{X}$ in the homotopy pullback square
  \[
    \begin{tikzcd}
      \widetilde{L}_i \arrow[r, "\widetilde{f}_i"] \arrow[d, "q_i"'] &
      \widetilde{X} \arrow[d, "p"] \\
      L_i \arrow[r, "f_i"'] & X
    \end{tikzcd}
  \]
The derived symplectification $\widetilde{X}$ carries a Liouville 1-form $\lambda_{\widetilde{X}}$, equipped with an equivalence $d_{dR}\lambda_{\widetilde{X}} \simeq \omega_{\widetilde{X}}$ in $\Acl(\widetilde{X}, n)$. The projection $p \colon \widetilde{X} \to X$ yields the equivalence $p^*\alpha \simeq \lambda_{\widetilde{X}}$. 
The commutativity of the square induces $q_i^* f_i^* \alpha \simeq \widetilde{f}_i^* \lambda_{\widetilde{X}}$. Base change of the isotropic structure $f_i^*\alpha \simeq 0$ along $q_i$ defines a null-homotopy $\widetilde{f}_i^*\lambda_{\widetilde{X}} \simeq 0$ via the diagram
  \[
    \begin{tikzcd}
      0 \arrow[r, "\sim"] \arrow[d, equal] & q_i^* f_i^* \alpha
      \arrow[d, "\sim"] \\
      0 \arrow[r, "\sim"'] & \widetilde{f}_i^* \lambda_{\widetilde{X}}
    \end{tikzcd}
  \]
  Applying the derived de Rham differential gives a $\Gm$-equivariant symplectic null-homotopy
  \begin{equation*}
    h_i \colon 0 \sim \widetilde{f}_i^*\omega_{\widetilde{X}} \quad \text{in} \quad \Acl(\widetilde{L}_i, n)^{\Gm}.
  \end{equation*}

  An isotropic structure $h_i$ on $\widetilde{f}_i$ is Lagrangian if the morphism $\Theta_{\widetilde{f}_i} \colon \TT_{\widetilde{L}_i} \to \LL_{\widetilde{L}_i/\widetilde{X}}[n-1]$ is a quasi-isomorphism \cite[Definition 2.8]{PTVV}. By Definition \ref{definition:legendrian structures}, the Legendrian structure on $L_i \to X$ is witnessed by the stable fiber sequence
  \begin{equation*}
    \cO_{L_i} \to \LL_{L_i/X} \otimes^{\mathbb{L}}_{\cO_{L_i}} f_i^*\cL[n-1] \to \TT_{L_i}.
  \end{equation*}
  We pull this sequence back along the principal $\Gm$-bundle $q_i \colon \widetilde{L}_i \to L_i$.
  By definition of the symplectification, the pullback $q_i^*f_i^*\cL$ is canonically trivialized to $\cO_{\widetilde{L}_i}$.
  Therefore, the pulled-back Legendrian sequence simplifies to the exact triangle
  \begin{equation} \label{eq:legendrian_pullback}
    \cO_{\widetilde{L}_i} \to q_i^*\LL_{L_i/X}[n-1] \to q_i^*\TT_{L_i}.
  \end{equation}
  Because $q_i$ is a principal $\Gm$-bundle, its derived Atiyah sequence gives the exact triangle
  \begin{equation} \label{eq:atiyah_bundle}
    \cO_{\widetilde{L}_i} \to \TT_{\widetilde{L}_i} \to q_i^*\TT_{L_i}.
  \end{equation}
  A map between the extensions \eqref{eq:atiyah_bundle} and \eqref{eq:legendrian_pullback} induces the equivalence $\TT_{\widetilde{L}_i} \simeq q_i^*\LL_{L_i/X}[n-1]$
  \[
    \begin{tikzcd}
      \cO_{\widetilde{L}_i} \arrow[r] \arrow[d, equal] &
      \TT_{\widetilde{L}_i} \arrow[r] \arrow[d, "\simeq"] &
      q_i^*\TT_{L_i} \arrow[d, equal] \\
      \cO_{\widetilde{L}_i} \arrow[r] & q_i^*\LL_{L_i/X}[n-1] \arrow[r]
      & q_i^*\TT_{L_i}
    \end{tikzcd}
  \]
  Since the square is a homotopy pullback, the relative cotangent complexes satisfy the equivalence $\LL_{\widetilde{L}_i/\widetilde{X}} \simeq q_i^* \LL_{L_i/X}$ \cite[Lemma 1.4.1.16(2)]{HAG-II}. Substituting this identification gives the equivalence $\Theta_{\widetilde{f}_i} \colon \TT_{\widetilde{L}_i} \xrightarrow{\sim} \LL_{\widetilde{L}_i/\widetilde{X}}[n-1]$ in $L_{qcoh}(\widetilde{L}_i)$. The lifts $\widetilde{f}_i$ are then $\Gm$-equivariant Lagrangian morphisms.
  
%\begin{figure}[htbp]
    %\centering
    %\begin{tikzpicture}[scale=.5]
    %    \draw[thick] (0,0) ellipse (2.5 and 0.8);
     %   \node at (2.6, -0.5) {\small $X$};

     %   \draw[thick] (-1.5, -0.2) to[out=30, in=150] (1.5, 0.1);
     %   \node at (1.8, 0) {\small $L_i$};

      %  \draw[->, thick] (-3, 1.5) -- (-3, 0.5) node[midway, left] {\small $p$};

     %   \draw[dashed, gray] (-2.5, 0) -- (-2.5, 2.5);
      %  \draw[dashed, gray] (2.5, 0) -- (2.5, 2.5);
     %   \draw[gray] (0, 2.5) ellipse (2.5 and 0.8);
      %  \node at (3, 2) {\small $\widetilde{X}$};

     %   \filldraw[fill=red!20, draw=red, thick, opacity=0.8] 
     %       (-1.5, -0.2) to[out=30, in=150] (1.5, 0.1) -- 
      %      (1.5, 2.6) to[in=30, out=150] (-1.5, 2.3) -- cycle;
      %  \node at (-0.1, 1.8) {\small $\widetilde{L}_i$};

      %  \draw[->, thick] (0.5, 0.425) -- (0.5, 1.5) node[midway, right] {\small $\Gm$};
      %  \draw[->, thick] (-0.5, 0.28) -- (-0.5, 1.3);
    %\end{tikzpicture}
    %\caption{\small The Lagrangian lift $\widetilde{L}_i$ inside the derived symplectification $\widetilde{X}$ over the Legendrian $L_i$ in the contact base $X$. The fibers are generated by the weight 1 $\Gm$-action modeling the Liouville vector field.}
%\end{figure}
\vspace{5pt}


\textbf{\textit{Step 2: The Transgression Loop in Equivariant Mapping Spaces.}}
  The derived fiber product of the lifts is given by the homotopy pullback $\widetilde{Z} := \widetilde{L}_1 \times_{\widetilde{X}}^h \widetilde{L}_2$ \cite[Section 2.2.2]{HAG-II}. Let $p_1 \colon \widetilde{Z} \to \widetilde{L}_1$ and $p_2 \colon \widetilde{Z} \to \widetilde{L}_2$ be the projections, and let $q \colon \widetilde{Z} \to \widetilde{X}$ denote the composition. Since homotopy pullbacks are computed within the stable $\infty$-category of $\Gm$-equivariant stacks, $\widetilde{Z}$ inherits a $\Gm$-action.

  Recall that the $\Gm$-action on a derived stack induces an auxiliary fiber grading on the derived de Rham complex relative to $B\Gm$ \cite[Section 2.1.1]{Calaque}. A form homogeneous of weight 1 under the $\Gm$-action carries the grading $\{1\}$. 
  
  In our case, the symplectic form $\omega_{\widetilde{X}}$ has weight 1 and hence resides in the equivariant mapping space ${\Acl}^{\{1\}}(\widetilde{X}, n)$. The $\Gm$-equivariant lifts $\widetilde{f}_1$ and $\widetilde{f}_2$ pull back the isotropic structures into the equivariant category, placing $h_1$ and $h_2$ in ${\Acl}^{\{1\}}(\widetilde{L}_i, n)^{\Gm}$.

  The homotopy pullback square in the $\infty$-category of $\Gm$-equivariant stacks provides a $\Gm$-equivariant 2-morphism $\gamma \colon \widetilde{f}_1 \circ p_1 \simeq \widetilde{f}_2 \circ p_2$. Composing the null-homotopies $p_1^*h_1$ and $p_2^*h_2$ via $\gamma$ defines a based loop $\Lambda_{\widetilde{Z}}$ at the zero vertex within the space of weight 1 equivariant forms
  \begin{equation*}
    \Lambda_{\widetilde{Z}} := p_1^*h_1 \cdot \gamma \cdot (p_2^*h_2)^{-1} \in \Omega_0 {\Acl}^{\{1\}}(\widetilde{Z}, n)^{\Gm}.
  \end{equation*}

  By \cite[Definition 1.8]{PTVV}, the space of closed 2-forms is defined as the mapping space $\mathrm{Map}_{\epsilon-\mathrm{dg}_k}(\K(2)[-n-2], \DR(\widetilde{Z}/\K))$. The based loop space functor $\Omega_0$ operates by pre-composing with the suspension functor $\Sigma$ on the domain $\K(2)[-n-2]$ \cite[Section 1.1.2]{LurieHA}. The suspension functor shifts the cohomological degree to $\K(2)[-n-1]$. It operates only on the domain complex and does not interact with the target complex $\DR(\widetilde{Z}/\K)$ or the $\Gm$-equivariant fiber grading $\{1\}$. This establishes the equivalence of mapping spaces
  \begin{equation*}
    \Omega_0 {\Acl}^{\{1\}}(\widetilde{Z}, n)^{\Gm} \simeq {\Acl}^{\{1\}}(\widetilde{Z}, n-1)^{\Gm}.
  \end{equation*}
The transgressed form $\omega_{\widetilde{Z}}$ above resides in this target space. It hence inherently carries the $\{1\}$ grading and remains a weight 1 eigenform.

  \vspace{5pt}

  \textbf{\textit{Step 3: Verification of Non-Degeneracy.}}
  By the main theorem of shifted symplectic intersections \cite[Theorem 2.9]{PTVV}, the derived fiber product of Lagrangian morphisms carries an $(n-1)$-shifted symplectic structure. Establishing this non-degeneracy equivariantly requires verifying that the induced morphism $\Theta_{\omega_{\widetilde{Z}}} \colon \TT_{\widetilde{Z}} \xrightarrow{\sim} \LL_{\widetilde{Z}}[n-1]$ is an equivalence. 
  
  The homotopy pullback square gives a stable fiber sequence of relative tangent complexes
  \begin{equation*}
    \TT_{\widetilde{Z}/\widetilde{L}_2} \to \TT_{\widetilde{Z}} \to p_2^*\TT_{\widetilde{L}_2}
  \end{equation*}
  in $L_{qcoh}(\widetilde{Z})$ \cite[Section 1.2]{PTVV}. Derived base change along the projection $p_2$ gives the equivalence $\TT_{\widetilde{Z}/\widetilde{L}_2} \simeq p_1^*\TT_{\widetilde{L}_1/\widetilde{X}}$. Substituting this into the sequence we obtain
  \begin{equation*}
    p_1^*\TT_{\widetilde{L}_1/\widetilde{X}} \to \TT_{\widetilde{Z}} \to p_2^*\TT_{\widetilde{L}_2}.
  \end{equation*}
  Conversely, we consider the relative cotangent fiber sequence for the projection $p_1$
  \begin{equation*}
    p_1^*\LL_{\widetilde{L}_1}[n-1] \to \LL_{\widetilde{Z}}[n-1] \to \LL_{\widetilde{Z}/\widetilde{L}_1}[n-1].
  \end{equation*}
  Again, derived base change yields $\LL_{\widetilde{Z}/\widetilde{L}_1} \simeq p_2^*\LL_{\widetilde{L}_2/\widetilde{X}}$. The symplectic form $\omega_{\widetilde{X}}$ and the Lagrangian isotropic structures $h_1, h_2$ induce a morphism between these two fiber sequences. The Lagrangian condition on $\widetilde{L}_2$ requires that the induced map $\Theta_2 \colon \TT_{\widetilde{L}_2} \xrightarrow{\sim} \LL_{\widetilde{L}_2/\widetilde{X}}[n-1]$ is an equivalence in $L_{qcoh}(\widetilde{L}_2)$ \cite[Definition 2.8]{PTVV}. 
  
  Rotating the exact triangle for the Lagrangian immersion $\widetilde{L}_1 \to \widetilde{X}$ and applying the $\cO$-linear duality functor along with a cohomological shift shows the isotropic structure $h_1$ equivalently induces an isomorphism $\Theta_1^\vee[n-1] \colon \TT_{\widetilde{L}_1/\widetilde{X}} \xrightarrow{\sim} \LL_{\widetilde{L}_1}[n-1]$.
  
  Pulling these back to $\widetilde{Z}$ constructs the commutative diagram of exact triangles
  \[
    \begin{tikzcd}
      p_1^*\TT_{\widetilde{L}_1/\widetilde{X}} \arrow[r] \arrow[d, "p_1^*\Theta_1^\vee{[n-1]} \simeq"'] &
      \TT_{\widetilde{Z}} \arrow[r] \arrow[d, "\Theta_{\omega_{\widetilde{Z}}}"] &
      p_2^*\TT_{\widetilde{L}_2} \arrow[d, "p_2^*\Theta_2 \simeq"'] \\
      p_1^*\LL_{\widetilde{L}_1}[n-1] \arrow[r] &
      \LL_{\widetilde{Z}}[n-1] \arrow[r] &
      p_2^*\LL_{\widetilde{L}_2/\widetilde{X}}[n-1]
    \end{tikzcd}
  \]
  in $L_{qcoh}(\widetilde{Z})$. This diagram commutes up to the homotopies provided by the paths $h_1, h_2$, and the 2-cell $\gamma$. By the pasting law for stable fiber sequences (Lemma \ref{lem:pasting_law}), since the outer vertical morphisms are equivalences of perfect complexes, the induced central morphism $\Theta_{\omega_{\widetilde{Z}}}$ is an equivalence. Therefore, $\omega_{\widetilde{Z}}$ is a {\em $\Gm$-equivariant $(n-1)$-shifted symplectic form.}

  \vspace{5pt}

  \textbf{\textit{Step 4: Descent via Base Change.}}
  Let $\pi \colon \widetilde{Z} \to Z$ be the projection to the contact stack quotient $Z := L_1 \times_X^h L_2$. To apply Theorem \ref{thm:A}, the projection $\pi$ must be a principal $\Gm$-bundle over $Z$. The homotopy pullback defining $\widetilde{Z}$ evaluates as
  \begin{equation*}
    \widetilde{Z} := \widetilde{L}_1 \times_{\widetilde{X}}^h \widetilde{L}_2 \simeq (L_1 \times_X^h \widetilde{X}) \times_{\widetilde{X}}^h (L_2 \times_X^h \widetilde{X}).
  \end{equation*}
  By the associativity and commutativity of limits in $\dArt_\K$, this simplifies to
  \begin{equation*}
    \widetilde{Z} \simeq (L_1 \times_X^h L_2) \times_X^h \widetilde{X} \simeq Z \times_X^h \widetilde{X}.
  \end{equation*}
  
  The descent structure is given by the Cartesian cube. Denote by $q_1 \colon \widetilde{L}_1 \to L_1$ and $q_2 \colon \widetilde{L}_2 \to L_2$ the principal $\Gm$-bundles associated with the lifts. The commutativity of the Lagrangian morphisms $\widetilde{f}_i$ with the torsor projections is given by the base-change 2-cell equivalences $\beta_1$ and $\beta_2$. Here, the top face encodes the Lagrangian intersection equipped with the 2-morphism $\gamma \colon \widetilde{f}_1 \circ p_1 \simeq \widetilde{f}_2 \circ p_2$.
  \[
    \begin{tikzcd}[row sep=2.5em, column sep=2.5em]
      \widetilde{Z} \arrow[rr, "{p_2}"] \arrow[dd, "{p_1}"'] \arrow[dr, "{\pi}" description] 
      & & \widetilde{L}_2 \arrow[dd, "{\widetilde{f}_2}" near start] \arrow[dr, "{q_2}" description] \\
      & Z \arrow[rr, "{\pi_2}" near start, crossing over] \arrow[dd, "{\pi_1}" near start, crossing over] 
      & & L_2 \arrow[dd, "{f_2}"] \\
      \widetilde{L}_1 \arrow[rr, "{\widetilde{f}_1}"' near start] \arrow[dr, "{q_1}" description] 
      & & \widetilde{X} \arrow[dr, "{p}" description] \\
      & L_1 \arrow[rr, "{f_1}"'] 
      & & X
      \arrow[from=1-1, to=3-3, phantom, "{\simeq \gamma}" description]
      \arrow[from=2-2, to=4-4, phantom, "{\simeq \bar{\gamma}}" description]
      \arrow[from=3-1, to=4-4, phantom, "{\simeq \beta_1}" description]
      \arrow[from=1-3, to=4-4, phantom, "{\simeq \beta_2}" description]
    \end{tikzcd}
  \]

  Since $\widetilde{X} \to X$ is a $\Gm$-torsor, and torsors are stable under derived base change \cite[Section 1.3.5]{HAG-II}, the projection $\pi \colon \widetilde{Z} \to Z$ is a principal $\Gm$-bundle over $Z$. The contact line bundle on the intersection $Z$ is defined as $\cL_Z := \pi_1^* f_1^* \cL \simeq \pi_2^* f_2^* \cL$, where $\pi_i \colon Z \to L_i$ are the projections. Since $\widetilde{Z} \simeq Z \times_X^h \widetilde{X}$, we conclude that $\widetilde{Z}$ is the principal $\Gm$-bundle associated with the line bundle $\cL_Z$ on $Z$.

  By \textbf{\textit{Step 2}} and \textbf{\textit{Step 3}}, $(\widetilde{Z}, \omega_{\widetilde{Z}})$ is an $(n-1)$-shifted symplectic stack equipped with a weight 1 $\Gm$-action. Applying Theorem \ref{thm:A}, the stack quotient $Z \simeq \varinjlim_{\Delta^{\mathrm{op}}} B_\bullet(\Gm, \widetilde{Z})$ descends the data to an $(n-1)$-shifted contact structure on $Z$. The descent of the top-face coherence $\gamma$ along the torsor projections and lateral base-change equivalences $\beta_1, \beta_2$ uniquely defines the Legendrian 2-morphism $\bar{\gamma} \colon f_1 \circ \pi_1 \simeq f_2 \circ \pi_2$. 
  
  The fundamental vector field of the $\Gm$-action pairs with $\omega_{\widetilde{Z}}$ to generate the 1-form $\alpha_{\widetilde{Z}}$ on the total space. Equivariant descent along the torsor projection $\pi$ yields the 1-form $\alpha_Z \in \pi_0\mathbb{R}\Gamma(Z, \LL_Z \otimes^{\mathbb{L}}_{\cO_Z} \cL_Z[n-1])$. 
  
  The contact distribution $\cK_Z$ is then the homotopy fiber $\mathrm{fib}(\TT_Z \to \cL_Z[n-1])$. The differential $d_{dR}$ commutes with the quotient map $\pi$. The underlying 2-form of $d_{dR}\alpha_Z$ pulls back to $\omega_{\widetilde{Z}}$. From \textbf{\textit{Step 3}}, $\omega_{\widetilde{Z}}$ induces an equivalence $\TT_{\widetilde{Z}} \xrightarrow{\sim} \LL_{\widetilde{Z}}[n-1]$. Descent of this equivariant non-degeneracy establishes the equivalence $\cK_Z \xrightarrow{\sim} \cK_Z^\vee \otimes^{\mathbb{L}}_{\cO_Z} \cL_Z[n-1]$ in $L_{qcoh}(Z)$.

  This concludes the proof of Theorem \ref{thm:Leg intersection}, and consequently, also proves Theorem \ref{thm:B}.
\end{proof}

\section{Applications} \label{sec:moduli}
\subsection{\texorpdfstring{$(-1)-$}-Shifted Contact Structures via  Derived Legendrian Intersections}

In the case of smooth stacks, an immediate corollary of Theorem \ref{thm:Leg intersection} arises as follows.
\begin{corollary}
    Let $X$ be a smooth contact $\K$-stack, 
and 
$L, L' \subset X$ two smooth Legendrian substacks.
Then the derived fiber product 
\(
L \times_X^h L'
\)
carries a canonical $(-1)$-shifted contact structure.
\end{corollary}

We will now examine a special and interesting case:
\begin{example}[Derived discriminant locus]\label{example: ddislocus}
Given any smooth scheme $L$ over $\K$, every regular function $f\in \cO(L)$ gives a\textit{ Legendrian embedding} of $L$ into the 1-jet space $(J^1(L), \xi_{jet})$ via the map \cite{Geiges} $$F:L\longrightarrow \mathbb{A}^1_\K \times T^*L, \quad p \mapsto (f(p),d_{dR}f_p).$$ Here,  we have the identification $J^1(L)\simeq \mathbb{A}^1_\K \times T^*L$ such that the 1-jet space $J^1(L)$ carries a canonical (0-shifted) contact structure $\xi_{jet}$. Note in particular that the zero function $f\equiv0$ corresponds to the zero section $j^10 \simeq L \subset T^*L \subset T^*L \times \mathbb{A}^1_\K \simeq J^1(L)$ of  $J^1(L).$

Classically, we define the (ordinary) \textit{discriminant locus} 
to be the image of the intersection of the graph $\Gamma_F$ and the zero-section of the 1-jet bundle $J^1(L),$ which usually fails to be a smooth scheme due to the expected transversality and degeneracy issues.

We then consider the \bfem{derived discriminant locus} $\Delta\mathrm{loc}(f)$ of $f$ defined as the derived fiber product
 \[\begin{tikzcd}
 \Delta\mathrm{loc}(f):= L \times_{J^1(L)}^h L   \arrow[r]   \arrow[d]             & L \arrow[d,"{0}"] \\
L \arrow[r, "F"] & (J^1(L), \xi_{jet}).                    
\end{tikzcd}\] 


Applying Theorem \ref{thm:Leg intersection} to the Legendrian morphisms $F: L\rightarrow (J^1(L), \xi_{jet}) \leftarrow L :{0}$, we conclude:
\begin{corollary}
The derived discriminant locus
$\Delta\mathrm{loc}(f)$ of a smooth function $f$ on a smooth scheme $L$ admits a $(-1)$-shifted contact structure.
   
\end{corollary}
 \end{example}   
 
\subsection{Projective Higgs Bundles and Legendrian Intersections}
Let $X$ be a smooth proper variety of dimension $d$ over $\K$ and let $G$ be a reductive algebraic group. The derived moduli stack of principal $G$-bundles $\mathrm{Bun}_G(X)$ is a derived Artin stack. 
By Serre duality, the cotangent complex $\LL_{\mathrm{Bun}_G(X)}$ at a point $P$ is equivalent to $R\Gamma(X, \mathrm{ad}(P)^* \otimes \omega_X)[d-1]$. To corepresent classical $\omega_X$-valued Higgs fields $\theta \in H^0(X, \mathrm{ad}(P)^* \otimes \omega_X)$ in cohomological degree 0, the cotangent stack must be shifted by $1-d$. Thus the derived moduli space of $G$-Higgs bundles on $X$ is the $(1-d)$-shifted cotangent stack $T^*[1-d]\mathrm{Bun}_G(X)$. Hence it carries a canonical $(1-d)$-shifted symplectic structure.

The multiplicative group $\Gm$ acts on $T^*[1-d]\mathrm{Bun}_G(X)$ by scaling the cotangent fibers. The fundamental weight 1 condition holds. Let $T^*[1-d]\mathrm{Bun}_G(X)^\circ$ denote the open substack where the Higgs field vanishes nowhere. Then Theorem \ref{thm:derived transversality} implies:

\begin{corollary} \label{cor:projective_higgs}
  The derived moduli stack of projective Higgs bundles
  \[
    P\mathrm{Higgs}_G(X) := [T^*[1-d]\mathrm{Bun}_G(X)^\circ / \Gm]
  \]
  admits a canonical $(1-d)$-shifted contact structure.
\end{corollary}

We restrict to the case $d=1$, where $C$ is a smooth proper curve. Substituting $d=1$ into Corollary \ref{cor:projective_higgs} yields the following specialization.

\begin{corollary} \label{cor:curve_higgs}
  The derived moduli stack of projective Higgs bundles $P\mathrm{Higgs}_G(C)$ admits a canonical $0$-shifted contact structure.
\end{corollary}

This ambient space serves as the geometric foundation to evaluate negative-shifted derived intersections.
Consider the global nilpotent cone $\mathcal{N} \subset T^*[0]\mathrm{Bun}_G(C)$, which is the central fiber of the Hitchin fibration. The nilpotency condition is invariant under the $\Gm$-scaling action. The projectivization of an irreducible component of the nilpotent cone, denoted $\mathbb{P}\mathcal{N}_i$, descends to a well-defined substack in $P\mathrm{Higgs}_G(C)$. By \cite[Theorem 6.9]{BenBassat}, the generic smooth part of the reduced nilpotent cone is an isotropic substack of the $0$-shifted symplectic Higgs moduli space. Since its dimension is half the dimension of the ambient space, it forms a Lagrangian substack. Therefore this inclusion defines a Legendrian morphism into the $0$-shifted contact stack.

\begin{corollary} \label{cor:nilpotent_intersection}
Let $\mathbb{P}\mathcal{N}_1$ and $\mathbb{P}\mathcal{N}_2$ be Legendrian substacks in $P\mathrm{Higgs}_G(C)$ associated with the projectivizations of two distinct irreducible components of the global nilpotent cone. Their derived fiber product
\[
    \mathbb{P}\mathcal{N}_1 \times_{P\mathrm{Higgs}_G(C)}^h \mathbb{P}\mathcal{N}_2
\]
carries a canonical $(-1)$-shifted contact structure.
\end{corollary}

\begin{proof}
The inclusions $f_i \colon \mathbb{P}\mathcal{N}_i \to P\mathrm{Higgs}_G(C)$ are Legendrian morphisms into a $0$-shifted contact stack. The homotopy pullback of their $\Gm$-equivariant Lagrangian lifts in the derived symplectification $T^*[0]\mathrm{Bun}_G(C)^\circ$ yields a $\Gm$-equivariant $(-1)$-shifted symplectic structure of weight 1. By Theorem 4.1, the descent of this structure along the principal $\Gm$-bundle quotient defines the $(-1)$-shifted contact structure on the intersection.
\end{proof}

\subsection{\texorpdfstring{$\ell$}{l}-adic Local Systems and the Contact Mapping Torus}

We thank the anonymous referee for suggesting this direction to promote a contact moduli space from arithmetic geometry.
Let $X$ be a smooth proper scheme of dimension $d$ defined over a finite field $\mathbb{F}_q$. Let $X_{\bar{\mathbb{F}}_q}$ denote its base change to the algebraic closure. The \bfem{derived moduli stack of framed continuous $\ell$-adic representations of the \'{e}tale fundamental group}, denoted $LocSys_{\ell,n}^{fr}(X_{\bar{\mathbb{F}}_q})$, is representable by a derived $\mathbb{Q}_\ell$-analytic space \cite[Theorem 1.1]{Antonio}. 

Recall that the tangent complex at a local system $L$ is computed by the \'{e}tale cohomology complex $R\Gamma_{et}(X_{\bar{\mathbb{F}}_q}, \mathrm{ad}(L))[1]$. Through \'{e}tale Poincar\'{e} duality and the trace map $\mathrm{tr} \colon H^{2d}_{et}(X_{\bar{\mathbb{F}}_q}, \mathbb{Q}_\ell(d)) \to \mathbb{Q}_\ell$, this moduli space inherits a $(2-2d)$-shifted symplectic structure \cite[Section 5]{Antonio}. Because we consider framed local systems, this symplectic form is exact, admitting a derived Liouville 1-form $\theta$ such that $d_{dR}\theta = \omega$ \cite[Section 2]{Calaque}.

The geometric Frobenius automorphism $F \in \mathrm{Gal}(\bar{\mathbb{F}}_q/\mathbb{F}_q)$ acts globally on $LocSys_{\ell,n}^{fr}(X_{\bar{\mathbb{F}}_q})$. The action of $F$ on the symplectic form is determined by its action on the \'{e}tale orientation class $\mathbb{Q}_\ell(d)$, which scales by the scalar $q^d$. Thus, the pullback of the geometric data scales uniformly by a constant:
\begin{equation*}
    F^* \omega = q^d \omega \quad \text{and} \quad F^*\theta = q^d \theta.
\end{equation*}

Because the arithmetic Frobenius generates a discrete group $\mathbb{Z}$, the direct quotient space preserves the even parity of the tangent complex, yielding a ``locally conformally symplectic stack''. To construct a shifted contact structure from this discrete arithmetic dilation, we pass to the \bfem{derived contact mapping torus} (the ``contactization").

Consider the product stack $Y = LocSys_{\ell,n}^{fr}(X_{\bar{\mathbb{F}}_q}) \times \mathbb{A}^1[2-2d]$, where $\tau$ is the coordinate on $\mathbb{A}^1[2-2d]$. The 1-form $\alpha = d_{dR}\tau + \theta$ defines a strict $(2-2d)$-shifted contact structure on $Y$, whose derived symplectification $\widetilde{Y} = Y \times \Gm$ carries the symplectic form $$\widetilde{\omega} = ds \wedge d_{dR}\tau + ds \wedge \theta + s \omega.$$
We then define a $\mathbb{Z}$-action on $Y$ where the generator acts by 
\begin{equation*}
    1 \cdot (x, \tau) = (F(x), q^d \tau).
\end{equation*}
Evaluating the pullback of the contact form under this $\mathbb{Z}$-action yields:
\begin{equation*}
    1^* \alpha = d_{dR}(q^d \tau) + F^*\theta = q^d d_{dR}\tau + q^d \theta = q^d \alpha.
\end{equation*}
Because the contact form $\alpha$ scales by the constant $q^d$, the contact distribution $\cK = \mathrm{fib}(\TT_Y \xrightarrow{\alpha^\vee} \cO_Y[2-2d])$ is strictly invariant under the $\mathbb{Z}$-action. 

Therefore, the geometric data descends to the quotient stack. The scaled form $\alpha$ maps to a section twisted by the flat line bundle $\cL$ determined by the character $1 \mapsto q^d$.

\begin{corollary} \label{cor: l-adic local sys}
    The derived contact mapping torus defined by the discrete arithmetic Frobenius scaling
    \begin{equation*}
        [(LocSys_{\ell,n}^{fr}(X_{\bar{\mathbb{F}}_q}) \times \mathbb{A}^1[2-2d]) / \mathbb{Z}]
    \end{equation*}
    inherits a canonical $(2-2d)$-shifted contact structure.
\end{corollary}

This geometric construction directly converts the arithmetic Frobenius scaling into a derived contact distribution. It extends to refined moduli problems in the categorical geometric Langlands program, linking the arithmetic geometry of local systems to the topology of shifted contact manifolds.

\subsection{Projective Coadjoint Stacks and Lie 2-Groups}

Connective structures on multiplicative gerbes define prequantizations of $BG$, yielding 2-shifted contact structures on the classifying stack of the associated Lie 2-group $\mathfrak{G}$ \cite[Section 3.4.2]{Tellez}. That construction relies on principal $B^2\Gm$-fibrations. We arrange a parallel geometry within our setup of derived symplectification via $\Gm$-torsors.
\vspace{1in}

Consider the classifying stack $BG$ of a reductive algebraic group $G$. The 2-shifted cotangent stack $T^*[2]BG \simeq \mathfrak{g}^\vee[1]/G$ carries a canonical 2-shifted symplectic structure  $\omega$ thanks to Theorem \ref{thm:calaque}. The multiplicative group $\Gm$ acts on the shifted fibers $\mathfrak{g}^\vee[1]$ by scaling. This action dilates the symplectic form $\omega$ with weight 1.

Let $T^*[2]BG^\circ$ denote the open substack corresponding to the complement of the zero section. By Theorem \ref{thm:derived transversality}, we then have the following result.

\begin{corollary} \label{cor: Lie 2-group}
  The derived stack quotient
  \begin{equation*}
    \mathbb{P}(T^*[2]BG) := [T^*[2]BG^\circ / \Gm]
  \end{equation*}
  inherits a canonical 2-shifted contact structure.
\end{corollary}

This projective stack provides an algebraic prequantization of $BG$, and hence Theorem \ref{thm:Leg intersection} applies to this space. Recall that for a closed subgroup $H \subset G$, the 2-shifted conormal stack $\mathcal{N}^*_{BH/BG}[2]$ defines a $\Gm$-invariant Lagrangian in $T^*[2]BG$. Combining the lifting argument introduced in the proof of Theorem \ref{thm:Leg intersection} with the contact reduction via Theorem \ref{thm:derived transversality}, its projectivization yields a Legendrian morphism $\mathbb{P}(\mathcal{N}^*_{BH/BG}[2]) \to \mathbb{P}(T^*[2]BG)$. 

\par From Theorem \ref{thm:Leg intersection}, we then conclude:
\begin{corollary}\label{cor: intersection of prequan}
    The derived fiber product of two such Legendrians $$\mathbb{P}(\mathcal{N}^*_{BH_1/BG}[2]) \to \mathbb{P}(T^*[2]BG) \leftarrow \mathbb{P}(\mathcal{N}^*_{BH_2/BG}[2])$$ associated with subgroups $H_1$ and $H_2$ carries a canonical 1-shifted contact structure.
\end{corollary}
    
 \section{Concluding Remarks}

In this paper, we have adapted the traditional transversality lemma from contact geometry to the framework of derived algebraic geometry and have shown the derived Legendrian intersection theorem, along with numerous applications, including the derived geometry of the discriminant loci of 1-jet bundles and certain moduli problems.

In brief, Theorem \ref{thm:derived transversality} establishes the categorical counterpart of the previously mentioned lemma for derived Artin stacks. Here, the central difference between the classical and derived settings can be summarized as follows: the transversality requirement (transverse hypersurface-local slice) in the classical setting arises because the global orbit space is often not a smooth, Hausdorff manifold. The derived stack quotient, on the other hand, incorporates these pathological orbits, thereby bypassing the local slice and achieving global contact reduction.

Theorem \ref{thm:Leg intersection} proves the derived Legendrian intersection theorem via base change and an $\infty$-categorical descent cube, along with a certain lifting property arising from the symplectic-contact dictionary. As shown in the proof of Theorem \ref{thm:Leg intersection}, this lifting property translates intersection problems in contact
moduli to equivariant intersection problems in symplectic moduli, and hence establishes the desired result. 

As a first application of the Legendrian intersection theorem, Theorem \ref{thm:Leg intersection}, we study in Example \ref{example: ddislocus} the derived geometry of the discriminant loci of 1-jet bundles and show that these loci possess a $(-1)$-shifted contact structure.

In the last section, we present further applications of Theorem \ref{thm:derived transversality} and Theorem \ref{thm:Leg intersection} to moduli theory, see  Corollary \ref{cor:projective_higgs}, Corollary \ref{cor: l-adic local sys}, Corollary \ref{cor: Lie 2-group}, and Corollary \ref{cor: intersection of prequan}, focusing on the moduli spaces of projective Higgs bundles, $\ell$-adic local systems, and Lie 2-groups.

\section*{Acknowledgements}
Kadri \.{I}lker Berktav posed the initial questions of formulating the transversality lemma and Legendrian intersections in the derived setup. He supervised the research, provided ideas for the proof of the Legendrian intersection theorem (Theorem \ref{thm:Leg intersection}), suggested directions for the moduli applications, and contributed the derived discriminant locus application (Example \ref{example: ddislocus}). 
Efe \.{I}zbudak formulated the statements of the derived transversality theorem (Theorem \ref{thm:derived transversality}) and the Legendrian intersection theorem (Theorem \ref{thm:Leg intersection}), provided the proofs, and constructed the applications to moduli problems (Section \ref{sec:moduli}).
The authors warmly thank the Higher Structures Research Group at METU for their support.

\begin{thebibliography}{99}

  \bibitem{LurieHTT}
  J. Lurie,
  \textit{Higher Topos Theory},
  Ann. Math. Stud., 170. Princeton University Press,
  Princeton, NJ, 2009.

  \bibitem{LurieHA}
  J. Lurie,
  \textit{Higher Algebra},
  Unpublished manuscript, available at author's website, 2017.

  \bibitem{Geiges}
  H. Geiges,
  \textit{An Introduction to Contact Topology},
  Cambridge University Press, 2008.

  \bibitem{HAG-II}
  B. To\"{e}n, G. Vezzosi,
  \textit{Homotopical algebraic geometry II: Geometric stacks and
  applications},
  Mem. Amer. Math. Soc. 193 (2008), no. 902, x+224 pp.

  \bibitem{PTVV}
  T. Pantev, B. To\"{e}n, M. Vaqui\'{e}, and G. Vezzosi,
  \textit{Shifted symplectic structures},
  Publications math\'{e}matiques de l'IH\'{E}S, 117:271--328, 2013.

  \bibitem{CPTVV}
  D. Calaque, T. Pantev, B. To\"{e}n, M. Vaqui\'{e}, and G. Vezzosi,
  \textit{Shifted Poisson Structures and deformation quantization},
  Journal of Topology, 10(2):483--584, 2017.

  \bibitem{Calaque}
  D. Calaque,
  \textit{Shifted cotangent stacks are shifted symplectic},
  Annales de la Facult\'{e} des sciences de Toulouse:
  Math\'{e}matiques, 28(1):67--90, 2019.

  \bibitem{Berktav1}
  K. \.{I}. Berktav,
  \textit{Shifted Contact Structures and Their Local Theory}, Ann. Fac.
  Sci. Toulouse, Math., Serie 6, Vol. 33(4): 1019-1057, 2024.

  \bibitem{Berktav2}
  K. \.{I}. Berktav,
  \textit{On shifted contact derived Artin stacks},
  Higher Structures 9(2):103--135, 2025.

  \bibitem{Berktav3}
  K. \.{I}. Berktav,
  \textit{Legendrian Structures in Derived Geometry},
  2025.

  \bibitem{ToenDAG}
  B. To\"{e}n,
  \textit{Derived Algebraic Geometry},
  EMS Surveys in Mathematical Sciences, 1(2):153--240, 2014.

  \bibitem{Safronov2023}
    P. Safronov.
    \newblock Shifted geometric quantization. 
    \newblock {\em Journal of Geometry and Physics}, 194, 104992,
    2023. \Arxiv{2011.05730}.

  \bibitem{Iwanari}
    I. Iwanari,
    \textit{Tannakization in derived algebraic geometry},
    Journal of K-theory, 14(3) 642--700, 2014, doi:10.1017/is014008019jkt278.

  \bibitem {TV2}  Toën, B. Vezzosi, G. 2004.  “From HAG to DAG: derived moduli stacks. In Axiomatic, enriched and motivic homotopy theory”,  Springer Netherlands, 173-216.

  \bibitem{Waterhouse}
W. C. Waterhouse, \textit{Introduction to Affine Group Schemes}, Graduate Texts in Mathematics, vol. 66, Springer-Verlag, 1979.

  \bibitem{Hitchin} N. J. Hitchin, \textit{The self-duality equations on a Riemann surface}, Proceedings of the London Mathematical Society, s3-55(1) 59--126, 1987.
  
  \bibitem{Tellez} R. T\'ellez Dom\'inguez, \textit{Moduli of Higher Connections and Holomorphic 2-Bundles}, PhD Thesis, Universidad Aut\'onoma de Madrid, 2025.

  \bibitem{Milne} J. S. Milne, \textit{\'{E}tale Cohomology}, Princeton Mathematical Series, 33. Princeton University Press, 1980.
  
  \bibitem{Antonio} J. Ant\'onio, \textit{Moduli of $\ell$-adic pro-\'{e}tale local systems for smooth non-proper schemes}, arXiv:1904.08001, 2019.
  
  \bibitem{Pridham} J. P. Pridham, \textit{Deforming $\ell$-adic representations of the fundamental group of a smooth variety}, Journal of Algebraic Geometry, 15:415--442, 2006.
  
  \bibitem{ArinkinRestricted} D. Arinkin, D. Gaitsgory, D. Kazhdan, S. Raskin, N. Rozenblyum, Y. Varshavsky, \textit{The stack of local systems with restricted variation and geometric Langlands theory with nilpotent singular support}, arXiv:2010.01906, 2020.
  
  \bibitem{ArinkinSingular} D. Arinkin, D. Gaitsgory, \textit{Singular support of coherent sheaves and the geometric Langlands conjecture}, arXiv:1201.6343, 2012.
  
  \bibitem{Zhu} X. Zhu, \textit{Coherent sheaves on the stack of Langlands parameters}, Proceedings of Symposia in Pure Mathematics, 2020.
  
  \bibitem{JiangAzumaya} Y. Jiang, \textit{On the construction of
  moduli stack of projective Higgs bundles over surfaces},
  arXiv:1911.00250, 2019.
  
  \bibitem{JiangKool} Y. Jiang, M. Kool, \textit{Twisted sheaves and
  $SU(r)/\mathbb{Z}_r$ Vafa-Witten theory}, arXiv:2006.10368, 2020.
  
  \bibitem{BenBassat} O. Ben-Bassat, S. Das, T. Pantev,
  \textit{Moduli stacks of Higgs bundles on stable curves}, Moduli 3, e7, 1-44, 2026.



\end{thebibliography}

\end{document}
